\documentclass[12pt,a4paper]{ctexbook}
\usepackage{geometry}
\geometry{margin=2.2cm}
\usepackage{setspace}
\setstretch{1.35}
\usepackage{hyperref}
\title{楚辞}
\author{古典文献整理}
\date{}
\begin{document}
\maketitle
\tableofcontents
\newpage
\section{离骚}
\textbf{离骚　屈原}\par\vspace{0.5em}
帝高阳之苗裔兮，朕皇考曰伯庸。\par
摄提贞于孟陬兮，惟庚寅吾以降。\par
皇览揆余初度兮，肇锡余以嘉名：\par
名余曰正则兮，字余曰灵均。\par
纷吾既有此内美兮，又重之以修能。\par
扈江离与辟芷兮，纫秋兰以为佩。\par
汨余若将不及兮，恐年岁之不吾与。\par
朝搴阰之木兰兮，夕揽洲之宿莽。\par
日月忽其不淹兮，春与秋其代序。\par
惟草木之零落兮，恐美人之迟暮。\par
不抚壮而弃秽兮，何不改乎此度?\par
乘骐骥以驰骋兮，来吾道夫先路！\par
昔三后之纯粹兮，固众芳之所在。\par
杂申椒与菌桂兮，岂惟纫夫蕙茝！\par
彼尧、舜之耿介兮，既遵道而得路。\par
何桀纣之猖披兮，夫惟捷径以窘步。\par
惟夫党人之偷乐兮，路幽昧以险隘。\par
岂余身之殚殃兮，恐皇舆之败绩！\par
忽奔走以先后兮，及前王之踵武。\par
荃不查余之中情兮，反信谗而齌怒。\par
余固知謇謇之为患兮，忍而不能舍也。\par
指九天以为正兮，夫惟灵修之故也。\par
曰黄昏以为期兮，羌中道而改路！\par
初既与余成言兮，后悔遁而有他。\par
余既不难夫离别兮，伤灵修之数化。\par
余既滋兰之九畹兮，又树蕙之百亩。\par
畦留夷与揭车兮，杂杜衡与芳芷。\par
冀枝叶之峻茂兮，愿俟时乎吾将刈。\par
虽萎绝其亦何伤兮，哀众芳之芜秽。\par
众皆竞进以贪婪兮，凭不厌乎求索。\par
羌内恕己以量人兮，各兴心而嫉妒。\par
忽驰骛以追逐兮，非余心之所急。\par
老冉冉其将至兮，恐修名之不立。\par
朝饮木兰之坠露兮，夕餐秋菊之落英。\par
苟余情其信姱以练要兮，长顑颔亦何伤。\par
掔木根以结茞兮，贯薜荔之落蕊。\par
矫菌桂以纫蕙兮，索胡绳之纚纚。\par
謇吾法夫前修兮，非世俗之所服。\par
虽不周于今之人兮，愿依彭咸之遗则。\par
长太息以掩涕兮，哀民生之多艰。\par
余虽好修姱以鞿羁兮，謇朝谇而夕替。\par
既替余以蕙纕兮，又申之以揽茝。\par
亦余心之所善兮，虽九死其犹未悔。\par
怨灵修之浩荡兮，终不察夫民心。\par
众女嫉余之蛾眉兮，谣诼谓余以善淫。\par
固时俗之工巧兮，偭规矩而改错。\par
背绳墨以追曲兮，竞周容以为度。\par
忳郁邑余侘傺兮，吾独穷困乎此时也。\par
宁溘死以流亡兮，余不忍为此态也。\par
鸷鸟之不群兮，自前世而固然。\par
何方圜之能周兮，夫孰异道而相安？\par
屈心而抑志兮，忍尤而攘诟。\par
伏清白以死直兮，固前圣之所厚。\par
悔相道之不察兮，延伫乎吾将反。\par
回朕车以复路兮，及行迷之未远。\par
步余马于兰皋兮，驰椒丘且焉止息。\par
进不入以离尤兮，退将复修吾初服。\par
制芰荷以为衣兮，集芙蓉以为裳。\par
不吾知其亦已兮，苟余情其信芳。\par
高余冠之岌岌兮，长余佩之陆离。\par
芳与泽其杂糅兮，唯昭质其犹未亏。\par
忽反顾以游目兮，将往观乎四荒。\par
佩缤纷其繁饰兮，芳菲菲其弥章。\par
民生各有所乐兮，余独好修以为常。\par
虽体解吾犹未变兮，岂余心之可惩。\par
女嬃之婵媛兮，申申其詈予，曰：\par
鲧婞直以亡身兮，终然夭乎羽之野。\par
汝何博謇而好修兮，纷独有此姱节？\par
薋菉葹以盈室兮，判独离而不服。\par
众不可户说兮，孰云察余之中情？\par
世并举而好朋兮，夫何茕独而不予听？\par
依前圣以节中兮，喟凭心而历兹。\par
济沅、湘以南征兮，就重华而敶词：\par
启《九辩》与《九歌》兮，夏康娱以自纵。\par
不顾难以图后兮，五子用失乎家衖。\par
羿淫游以佚畋兮，又好射夫封狐。\par
固乱流其鲜终兮，浞又贪夫厥家。\par
浇身被服强圉兮，纵欲而不忍。\par
日康娱而自忘兮，厥首用夫颠陨。\par
夏桀之常违兮，乃遂焉而逢殃。\par
后辛之菹醢兮，殷宗用而不长。\par
汤、禹俨而祗敬兮，周论道而莫差。\par
举贤才而授能兮，循绳墨而不颇。\par
皇天无私阿兮，览民德焉错辅。\par
夫维圣哲以茂行兮，苟得用此下土。\par
瞻前而顾后兮，相观民之计极。\par
夫孰非义而可用兮？\par
孰非善而可服？\par
阽余身而危死兮，览余初其犹未悔。\par
不量凿而正枘兮，固前修以菹醢。\par
曾歔欷余郁邑兮，哀朕时之不当。\par
揽茹蕙以掩涕兮，沾余襟之浪浪。\par
跪敷衽以陈辞兮，耿吾既得此中正。\par
驷玉虬以桀鹥兮，溘埃风余上征。\par
朝发轫于苍梧兮，夕余至乎县圃。\par
欲少留此灵琐兮，日忽忽其将暮。\par
吾令羲和弭节兮，望崦嵫而勿迫。\par
路漫漫其修远兮，吾将上下而求索。\par
饮余马于咸池兮，总余辔乎扶桑。\par
折若木以拂日兮，聊逍遥以相羊。\par
前望舒使先驱兮，后飞廉使奔属。\par
鸾皇为余先戒兮，雷师告余以未具。\par
吾令凤鸟飞腾兮，继之以日夜。\par
飘风屯其相离兮，帅云霓而来御。\par
纷总总其离合兮，斑陆离其上下。\par
吾令帝阍开关兮，倚阊阖而望予。\par
时暧暧其将罢兮，结幽兰而延伫。\par
世溷浊而不分兮，好蔽美而嫉妒。\par
朝吾将济于白水兮，登阆风而绁马。\par
忽反顾以流涕兮，哀高丘之无女。\par
溘吾游此春宫兮，折琼枝以继佩。\par
及荣华之未落兮，相下女之可诒。\par
吾令丰隆乘云兮，求宓妃之所在。\par
解佩纕以结言兮，吾令謇修以为理。\par
纷总总其离合兮，忽纬繣其难迁。\par
夕归次于穷石兮，朝濯发乎洧盘。\par
保厥美以骄傲兮，日康娱以淫游。\par
虽信美而无礼兮，来违弃而改求。\par
览相观于四极兮，周流乎天余乃下。\par
望瑶台之偃蹇兮，见有娀之佚女。\par
吾令鸩为媒兮，鸩告余以不好。\par
雄鸠之鸣逝兮，余犹恶其佻巧。\par
心犹豫而狐疑兮，欲自适而不可。\par
凤皇既受诒兮，恐高辛之先我。\par
欲远集而无所止兮，聊浮游以逍遥。\par
及少康之未家兮，留有虞之二姚。\par
理弱而媒拙兮，恐导言之不固。\par
世溷浊而嫉贤兮，好蔽美而称恶。\par
闺中既以邃远兮，哲王又不寤。\par
怀朕情而不发兮，余焉能忍而与此终古？\par
索琼茅以筳篿兮，命灵氛为余占之。\par
曰：两美其必合兮，孰信修而慕之？\par
思九州之博大兮，岂惟是其有女？」\par
曰：勉远逝而无狐疑兮，孰求美而释女？\par
何所独无芳草兮，尔何怀乎故宇？」\par
世幽昧以昡曜兮，孰云察余之善恶？\par
民好恶其不同兮，惟此党人其独异！\par
户服艾以盈要兮，谓幽兰其不可佩。\par
览察草木其犹未得兮，岂珵美之能当？\par
苏粪壤以充祎兮，谓申椒其不芳。\par
欲从灵氛之吉占兮，心犹豫而狐疑。\par
巫咸将夕降兮，怀椒糈而要之。\par
百神翳其备降兮，九疑缤其并迎。\par
皇剡剡其扬灵兮，告余以吉故。\par
曰：勉升降以上下兮，求矩矱之所同。\par
汤、禹俨而求合兮，挚、咎繇而能调。\par
苟中情其好修兮，又何必用夫行媒？\par
说操筑于傅岩兮，武丁用而不疑。\par
吕望之鼓刀兮，遭周文而得举。\par
宁戚之讴歌兮，齐桓闻以该辅。\par
及年岁之未晏兮，时亦犹其未央。\par
恐鹈鴃之先鸣兮，使夫百草为之不芳。\par
何琼佩之偃蹇兮，众薆然而蔽之。\par
惟此党人之不谅兮，恐嫉妒而折之。\par
时缤纷其变易兮，又何可以淹留？\par
兰芷变而不芳兮，荃蕙化而为茅。\par
何昔日之芳草兮，今直为此萧艾也？\par
岂其有他故兮，莫好修之害也！\par
余以兰为可恃兮，羌无实而容长。\par
委厥美以从俗兮，苟得列乎众芳。\par
椒专佞以慢慆兮，樧又欲充夫佩帏。\par
既干进而务入兮，又何芳之能祗？\par
固时俗之流从兮，又孰能无变化？\par
览椒兰其若兹兮，又况揭车与江离？\par
惟兹佩之可贵兮，委厥美而历兹。\par
芳菲菲而难亏兮，芬至今犹未沬。\par
和调度以自娱兮，聊浮游而求女。\par
及余饰之方壮兮，周流观乎上下。\par
灵氛既告余以吉占兮，历吉日乎吾将行。\par
折琼枝以为羞兮，精琼爢以为粻。\par
为余驾飞龙兮，杂瑶象以为车。\par
何离心之可同兮？\par
吾将远逝以自疏。\par
邅吾道夫昆仑兮，路修远以周流。\par
扬云霓之晻蔼兮，鸣玉鸾之啾啾。\par
朝发轫于天津兮，夕余至乎西极。\par
凤皇翼其承旗兮，高翱翔之翼翼。\par
忽吾行此流沙兮，遵赤水而容与。\par
麾蛟龙使梁津兮，诏西皇使涉予。\par
路修远以多艰兮，腾众车使径待。\par
路不周以左转兮，指西海以为期。\par
屯余车其千乘兮，齐玉轪而并驰。\par
驾八龙之婉婉兮，载云旗之委蛇。\par
抑志而弭节兮，神高驰之邈邈。\par
奏《九歌》而舞《韶》兮，聊假日以偷乐。\par
陟升皇之赫戏兮，忽临睨夫旧乡。\par
仆夫悲余马怀兮，蜷局顾而不行。\par
乱曰：已矣哉！国无人莫我知兮，又何怀乎故都！\par
既莫足与为美政兮，吾将从彭咸之所居！\par

\section{东皇太一}
\textbf{九歌　屈原}\par\vspace{0.5em}
吉日兮辰良，穆将愉兮上皇\par
抚长剑兮玉珥，璆锵鸣兮琳琅\par
瑶席兮玉瑱，盍将把兮琼芳\par
蕙肴蒸兮兰藉，奠桂酒兮椒浆\par
扬枹兮拊鼓，疏缓节兮安歌\par
陈竽瑟兮浩倡\par
灵偃蹇兮姣服，芳菲菲兮满堂\par
五音纷兮繁会，君欣欣兮乐康。\par

\section{云中君}
\textbf{九歌　屈原}\par\vspace{0.5em}
浴兰汤兮沐芳，华采衣兮若英\par
灵连蜷兮既留，烂昭昭兮未央\par
謇将憺兮寿宫，与日月兮齐光\par
龙驾兮帝服，聊翱游兮周章\par
灵皇皇兮既降，猋远举兮云中\par
览冀洲兮有余，横四海兮焉穷\par
思夫君兮太息，极劳心兮忡忡\par

\section{湘君}
\textbf{九歌　屈原}\par\vspace{0.5em}
君不行兮夷犹，蹇谁留兮中洲\par
美要眇兮宜修，沛吾乘兮桂舟\par
令沅湘兮无波，使江水兮安流\par
望夫君兮未来，吹参差兮谁思\par
驾飞龙兮北征，邅吾道兮洞庭\par
薜荔柏兮蕙绸，荪桡兮兰旌\par
望涔阳兮极浦，横大江兮扬灵\par
扬灵兮未极，女婵媛兮为余太息\par
横流涕兮潺湲，隐思君兮陫侧\par
桂棹兮兰枻，斵冰兮积雪\par
采薜荔兮水中，搴芙蓉兮木末\par
心不同兮媒劳，恩不甚兮轻绝\par
石濑兮浅浅，飞龙兮翩翩\par
交不忠兮怨长，期不信兮告余以不闲\par
朝骋骛兮江皋，夕弭节兮北渚\par
鸟次兮屋上，水周兮堂下\par
捐余玦兮江中，遗余佩兮醴浦\par
采芳洲兮杜若，将以遗兮下女\par
时不可兮再得，聊逍遥兮容与。\par

\section{湘夫人}
\textbf{九歌　屈原}\par\vspace{0.5em}
帝子降兮北渚，目眇眇兮愁予\par
袅袅兮秋风，洞庭波兮木叶下\par
登白薠兮骋望，与佳期兮夕张\par
鸟何萃兮苹中，罾何为兮木上？\par
沅有茝兮醴有兰，思公子兮未敢言\par
荒忽兮远望，观流水兮潺湲\par
麋何食兮庭中，蛟何为兮水裔\par
朝驰余马兮江皋，夕济兮西澨\par
闻佳人兮召余，将腾驾兮偕逝\par
筑室兮水中，葺之兮荷盖\par
荪壁兮紫坛，播芳椒兮成堂\par
桂栋兮兰橑，辛夷楣兮药房\par
罔薜荔兮为帷，擗蕙櫋兮既张\par
白玉兮为镇，疏石兰兮为芳\par
芷葺兮荷屋，缭之兮杜衡\par
合百草兮实庭，建芳馨兮庑门\par
九嶷缤兮并迎，灵之来兮如云\par
捐余袂兮江中，遗余褋兮醴浦\par
搴汀洲兮杜若，将以遗兮远者\par
时不可兮骤得，聊逍遥兮容与！\par

\section{大司命}
\textbf{九歌　屈原}\par\vspace{0.5em}
广开兮天门，纷吾乘兮玄云\par
令飘风兮先驱，使涷雨兮洒尘\par
君回翔兮以下，逾空桑兮从女\par
纷总总兮九州，何寿夭兮在予\par
高飞兮安翔，乘清气兮御阴阳\par
吾与君兮齐速，导帝之兮九坑\par
灵衣兮被被，玉佩兮陆离\par
一阴兮一阳，众莫知兮余所为\par
折疏麻兮瑶华，将以遗兮离居\par
老冉冉兮既极，不寖近兮愈疏\par
乘龙兮辚辚，高驰兮冲天\par
结桂枝兮延伫，羌愈思兮愁人\par
愁人兮奈何，愿若今兮无亏\par
固人命兮有当，孰离合兮何为？\par

\section{少司命}
\textbf{九歌　屈原}\par\vspace{0.5em}
秋兰兮麋芜，罗生兮堂下\par
绿叶兮素华，芳菲菲兮袭予\par
夫人兮自有美子，荪何以兮愁苦\par
秋兰兮青青，绿叶兮紫茎\par
满堂兮美人，忽独与余兮目成\par
入不言兮出不辞，乘回风兮载云旗\par
悲莫悲兮生别离，乐莫乐兮新相知\par
荷衣兮蕙带，儵而来兮忽而逝\par
夕宿兮帝郊，君谁须兮云之际\par
与女沐兮咸池，曦女发兮阳之阿\par
望美人兮未来，临风怳兮浩歌\par
孔盖兮翠旌，登九天兮抚彗星\par
竦长剑兮拥幼艾，荪独宜兮为民正。\par

\section{东君}
\textbf{九歌　屈原}\par\vspace{0.5em}
暾将出兮东方，照吾槛兮扶桑\par
抚余马兮安驱，夜皎皎兮既明\par
驾龙輈兮乘雷，载云旗兮委蛇\par
长太息兮将上，心低徊兮顾怀\par
羌声色兮娱人，观者儋兮忘归\par
縆瑟兮交鼓，萧钟兮瑶簴\par
鸣篪兮吹竽，思灵保兮贤姱\par
翾飞兮翠曾，展诗兮会舞\par
应律兮合节，灵之来兮敝日\par
青云衣兮白霓裳，举长矢兮射天狼\par
操余弧兮反沦降，援北斗兮酌桂浆\par
撰余辔兮高驰翔，杳冥冥兮以东行。\par

\section{河伯}
\textbf{九歌　屈原}\par\vspace{0.5em}
与女游兮九河，冲风起兮水扬波\par
乘水车兮荷盖，驾两龙兮骖螭\par
登昆仑兮四望，心飞扬兮浩荡\par
日将暮兮怅忘归，惟极浦兮寤怀\par
鱼鳞屋兮龙堂，紫贝阙兮珠宫\par
灵何惟兮水中\par
乘白鼋兮逐文鱼，与女游兮河之渚\par
流澌纷兮将来下\par
子交手兮东行，送美人兮南浦\par
波滔滔兮来迎，鱼鳞鳞兮媵予。\par

\section{山鬼}
\textbf{九歌　屈原}\par\vspace{0.5em}
若有人兮山之阿，被薜荔兮带女萝\par
既含睇兮又宜笑，子慕予兮善窈窕\par
乘赤豹兮从文狸，辛夷车兮结桂旗\par
被石兰兮带杜衡，折芬馨兮遗所思\par
余处幽篁兮终不见天，路险难兮独后来\par
表独立兮山之上，云容容兮而在下\par
杳冥冥兮羌昼晦，东风飘兮神灵雨\par
留灵修兮憺忘归，岁既晏兮孰华予\par
采三秀兮于山间，石磊磊兮葛蔓蔓\par
怨公子兮怅忘归，君思我兮不得闲\par
山中人兮芳杜若，饮石泉兮荫松柏\par
君思我兮然疑作\par
雷填填兮雨冥冥，猿啾啾兮狖夜鸣\par
风飒飒兮木萧萧，思公子兮徒离忧。\par

\section{国殇}
\textbf{九歌　屈原}\par\vspace{0.5em}
操吴戈兮被犀甲，车错毂兮短兵接\par
旌蔽日兮敌若云，矢交坠兮士争先\par
凌余阵兮躐余行，左骖殪兮右刃伤\par
霾两轮兮絷四马，援玉枹兮击鸣鼓\par
天时怼兮威灵怒，严杀尽兮弃原野\par
出不入兮往不反，平原忽兮路超远\par
带长剑兮挟秦弓，首身离兮心不惩\par
诚既勇兮又以武，终刚强兮不可凌\par
身既死兮神以灵，魂魄毅兮为鬼雄。\par

\section{礼魂}
\textbf{九歌　屈原}\par\vspace{0.5em}
成礼兮会鼓，\par
传芭兮代舞\par
姱女倡兮容与\par
春兰兮秋菊，\par
长无绝兮终古。\par

\section{天问}
\textbf{天问　屈原}\par\vspace{0.5em}
曰：\par
遂古之初，谁传道之？\par
上下未形，何由考之？\par
冥昭瞢暗，谁能极之？\par
冯翼惟象，何以识之？\par
明明暗暗，惟时何为？\par
阴阳三合，何本何化？\par
圜则九重，孰营度之？\par
惟兹何功，孰初作之？\par
斡维焉系，天极焉加？\par
八柱何当，东南何亏？\par
九天之际，安放安属？\par
隅隈多有，谁知其数？\par
天何所沓？\par
十二焉分？\par
日月安属？\par
列星安陈？\par
出自汤谷，次于蒙汜。\par
自明及晦，所行几里？\par
夜光何德，死则又育？\par
厥利维何，而顾菟在腹？\par
女岐无合，夫焉取九子？\par
伯强何处？\par
惠气安在？\par
何阖而晦？\par
何开而明？\par
角宿未旦，曜灵安藏？\par
不任汩鸿，师何以尚之？\par
佥曰\par
“何忧，何不课而行之？”\par
鸱龟曳衔，鲧何听焉？\par
顺欲成功，帝何刑焉？\par
永遏在羽山，夫何三年不施？\par
伯禹愎鲧，夫何以变化？\par
纂就前绪，遂成考功。\par
何续初继业，而厥谋不同？\par
洪泉极深，何以窴之？\par
地方九则，何以坟之？\par
河海应龙？\par
何尽何历？\par
鲧何所营？\par
禹何所成？\par
康回冯怒，墬何故以东南倾？\par
九州安错？\par
川谷何洿？\par
东流不溢，孰知其故？\par
东西南北，其修孰多？\par
南北顺椭，其衍几何？\par
昆仑悬圃，其尻安在？\par
增城九重，其高几里？\par
四方之门，其谁从焉？\par
西北辟启，何气通焉？\par
日安不到？\par
烛龙何照？\par
羲和之未扬，若华何光？\par
何所冬暖？\par
何所夏寒？\par
焉有石林？\par
何兽能言？\par
焉有虬龙，负熊以游？\par
雄虺九首，鯈忽焉在？\par
何所不死？\par
长人何守？\par
靡蓱九衢，枲华安居？\par
灵蛇吞象，厥大何如？\par
黑水玄趾，三危安在？\par
延年不死，寿何所止？\par
鲮鱼何所？\par
鬿堆焉处？\par
羿焉彃日？\par
乌焉解羽？\par
禹之力献功，降省下土四方。\par
焉得彼嵞山女，而通之於台桑？\par
闵妃匹合，厥身是继。\par
胡维嗜不同味，而快鼌饱？\par
启代益作后，卒然离蠥。\par
何启惟忧，而能拘是达？\par
皆归射鞫，而无害厥躬。\par
何后益作革，而禹播降？\par
启棘宾商，《九辨》《九歌》。\par
何勤子屠母，而死分竟地？\par
帝降夷羿，革孽夏民。\par
胡射夫河伯，而妻彼雒嫔？\par
冯珧利决，封豨是射。\par
何献蒸肉之膏，而后帝不若？\par
浞娶纯狐，眩妻爰谋。\par
何羿之射革，而交吞揆之？\par
阻穷西征，岩何越焉？\par
化而为黄熊，巫何活焉？\par
咸播秬黍，莆雚是营。\par
何由并投，而鲧疾修盈？\par
白蜺婴茀，胡为此堂？\par
安得夫良药，不能固臧？\par
天式从横，阳离爰死。\par
大鸟何鸣，夫焉丧厥体？\par
蓱号起雨，何以兴之？\par
撰体协胁，鹿何膺之？\par
鳌戴山抃，何以安之？\par
释舟陵行，何之迁之？\par
惟浇在户，何求于嫂？\par
何少康逐犬，而颠陨厥首？\par
女歧缝裳，而馆同爰止。\par
何颠易厥首，而亲以逢殆？\par
汤谋易旅，何以厚之？\par
覆舟斟寻，何道取之？\par
桀伐蒙山，何所得焉？\par
妺嬉何肆，汤何殛焉？\par
舜闵在家，父何以鳏？\par
尧不姚告，二女何亲？\par
厥萌在初，何所亿焉？\par
璜台十成，谁所极焉？\par
登立为帝，孰道尚之？\par
女娲有体，孰制匠之？\par
舜服厥弟，终然为害。\par
何肆犬豕，而厥身不危败？\par
吴获迄古，南岳是止。\par
孰期去斯，得两男子？\par
缘鹄饰玉，后帝是飨。\par
何承谋夏桀，终以灭丧？\par
帝乃降观，下逢伊挚。\par
何条放致罚，而黎服大说？\par
简狄在台，喾何宜？\par
玄鸟致贻，女何喜？\par
该秉季德，厥父是臧。\par
胡终弊于有扈，牧夫牛羊？\par
干协时舞，何以怀之？\par
平胁曼肤，何以肥之？\par
有扈牧竖，云何而逢？\par
击床先出，其命何从？\par
恒秉季德，焉得夫朴牛？\par
何往营班禄，不但还来？\par
昏微循迹，有狄不宁。\par
何繁鸟萃棘，负子肆情？\par
眩弟并淫，危害厥兄。\par
何变化以作诈，而后嗣逢长？\par
成汤东巡，有莘爰极。\par
何乞彼小臣，而吉妃是得？\par
水滨之木，得彼小子。\par
夫何恶之，媵有莘之妇？\par
汤出重泉，夫何辠尤？\par
不胜心伐帝，夫谁使挑之？\par
会朝争盟，何践吾期？\par
苍鸟群飞，孰使萃之？\par
列击纣躬，叔旦不嘉。\par
何亲揆发足，周之命以咨嗟？\par
授殷天下，其位安施？\par
反成乃亡，其罪伊何？\par
争遣伐器，何以行之？\par
并驱击翼，何以将之？\par
昭后成游，南土爰底。\par
厥利惟何，逢彼白雉？\par
穆王巧梅，夫何为周流？\par
环理天下，夫何索求？\par
妖夫曳炫，何号于市？\par
周幽谁诛？\par
焉得夫褒姒？\par
天命反侧，何罚何佑？\par
齐桓九会，卒然身杀。\par
彼王纣之躬，孰使乱惑？\par
何恶辅弼，谗谄是服？\par
比干何逆，而抑沈之？\par
雷开阿顺，而赐封之？\par
何圣人之一德，卒其异方？\par
梅伯受醢，箕子详狂？\par
稷维元子，帝何竺之？\par
投之于冰上，鸟何燠之？\par
何冯弓挟矢，殊能将之？\par
既惊帝切激，何逢长之？\par
伯昌号衰，秉鞭作牧。\par
何令彻彼岐社，命有殷国？\par
迁藏就岐，何能依？\par
殷有惑妇，何所讥？\par
受赐兹醢，西伯上告。\par
何亲就上帝罚，殷之命以不救？\par
师望在肆，昌何识？\par
鼓刀扬声，后何喜？\par
武发杀殷，何所悒？\par
载尸集战，何所急？\par
伯林雉经，维其何故？\par
何感天抑墬，夫谁畏惧？\par
皇天集命，惟何戒之？\par
受礼天下，又使至代之？\par
初汤臣挚，后兹承辅。\par
何卒官汤，尊食宗绪？\par
勋阖梦生，少离散亡。\par
何壮武历，能流厥严？\par
彭铿斟雉，帝何飨？\par
受寿永多，夫何久长？\par
中央共牧，后何怒？\par
蜂蛾微命，力何固？\par
惊女采薇，鹿何佑？\par
北至回水，萃何喜？\par
兄有噬犬，弟何欲？\par
易之以百两，卒无禄？\par
薄暮雷电，归何忧？\par
厥严不奉，帝何求？\par
伏匿穴处，爰何云？\par
荆勋作师，夫何长？\par
悟过改更，我又何言？\par
吴光争国，久余是胜。\par
何环穿自闾社丘陵，爰出子文？\par
吾告堵敖以不长。\par
何试上自予，忠名弥彰？\par

\section{惜诵}
\textbf{九章　屈原}\par\vspace{0.5em}
惜诵以致愍兮，发愤以抒情。\par
所作忠而言之兮，指苍天以为正。\par
令五帝使折中兮，戒六神与向服。\par
俾山川以备御兮，命咎繇使听直。\par
竭忠诚而事君兮，反离群而赘肬。\par
忘儇媚以背众兮，待明君其知之。\par
言与行其可迹兮，情与貌其不变。\par
故相臣莫若君兮，所以证之不远。\par
吾谊先君而后身兮，羌众人之所仇也。\par
专惟君而无他兮，又众兆之所雠也。\par
壹心而不豫兮，羌无可保也。\par
疾亲君而无他兮，有招祸之道也。\par
思君其莫我忠兮，忽忘身之贱贫。\par
事君而不贰兮，迷不知宠之门。\par
患何罪以遇罚兮，亦非余之所志也。\par
行不群以巅越兮，又众兆之所咍也。\par
纷逢尤以离谤兮，謇不可释也。\par
情沉抑而不达兮，又蔽而莫之白也。\par
心郁邑余侘傺兮，又莫察余之中情。\par
固烦言不可结而诒兮，愿陈志而无路。\par
退静默而莫余知兮，进号呼又莫吾闻。\par
申侘傺之烦惑兮，中闷瞀之忳忳。\par
昔余梦登天兮，魂中道而无杭。\par
吾使厉神占之兮，曰有志极而无旁。\par
终危独以离异兮，曰君可思而不可恃。\par
故众口其铄金兮，初若是而逢殆。\par
惩于羹者而吹齑兮，何不变此志也？\par
欲释阶而登天兮，犹有曩之态也。\par
众骇遽以离心兮，又何以为此伴也？\par
同极而异路兮，又何以为此援也？\par
晋申生之孝子兮，父信谗而不好。\par
行婞直而不豫兮，鲧功用而不就。\par
吾闻作忠以造怨兮，忽谓之过言。\par
九折臂而成医兮，吾至今而知其信然。\par
矰弋机而在上兮，罻罗张而在下。\par
设张辟以娱君兮，愿侧身而无所。\par
欲儃徊以干傺兮，恐重患而离尤。\par
欲高飞而远集兮，君罔谓汝何之？\par
欲横奔而失路兮，盖志坚而不忍。\par
背膺牉以交痛兮，心郁结而纡轸。\par
擣木兰以矫蕙兮，糳申椒以为粮。\par
播江离与滋菊兮，愿春日以为糗芳。\par
恐情质之不信兮，故重著以自明。\par
矫兹媚以私处兮，愿曾思而远身。\par

\section{涉江}
\textbf{九章　屈原}\par\vspace{0.5em}
余幼好此奇服兮，年既老而不衰。\par
带长铗之陆离兮，冠切云之崔嵬，被明月兮佩宝璐。\par
世混浊而莫余知兮，吾方高驰而不顾。\par
驾青虬兮骖白螭，吾与重华游兮瑶之圃。\par
登昆仑兮食玉英，与天地兮同寿，与日月兮同光。\par
哀南夷之莫吾知兮，旦余济乎江湘。\par
乘鄂渚而反顾兮，欸秋冬之绪风。\par
步余马兮山皋，邸余车兮方林。\par
乘舲船余上沅兮，齐吴榜以击汰。\par
船容与而不进兮，淹回水而疑滞。\par
朝发枉渚兮，夕宿辰阳。\par
苟余心其端直兮，虽僻远之何伤。\par
入溆浦余儃徊兮，迷不知吾所如。\par
深林杳以冥冥兮，猿狖之所居。\par
山峻高以蔽日兮，下幽晦以多雨。\par
霰雪纷其无垠兮，云霏霏而承宇。\par
哀吾生之无乐兮，幽独处乎山中。\par
吾不能变心而从俗兮，固将愁苦而终穷。\par
接舆髡首兮，桑扈臝行。\par
忠不必用兮，贤不必以。\par
伍子逢殃兮，比干菹醢。\par
与前世而皆然兮，吾又何怨乎今之人！\par
余将董道而不豫兮，固将重昏而终身！\par
乱曰：\par
鸾鸟凤皇，日以远兮。\par
燕雀乌鹊，巢堂坛兮。\par
露申辛夷，死林薄兮。\par
腥臊并御，芳不得薄兮。\par
阴阳易位，时不当兮。\par
怀信佗傺，忽乎吾将行兮！\par

\section{哀郢}
\textbf{九章　屈原}\par\vspace{0.5em}
皇天之不纯命兮，何百姓之震愆？\par
民离散而相失兮，方仲春而东迁。\par
去故乡而就远兮，遵江夏以流亡。\par
出国门而轸怀兮，甲之鼂吾以行。\par
发郢都而去闾兮，怊荒忽其焉极？\par
楫齐扬以容与兮，哀见君而不再得。\par
望长楸而太息兮，涕淫淫其若霰。\par
过夏首而西浮兮，顾龙门而不见。\par
心婵媛而伤怀兮，眇不知其所蹠。\par
顺风波以从流兮，焉洋洋而为客。\par
凌阳侯之汜滥兮，忽翱翔之焉薄。\par
心絓结而不解兮，思蹇产而不释。\par
将运舟而下浮兮，上洞庭而下江。\par
去终古之所居兮，今逍遥而来东。\par
羌灵魂之欲归兮，何须臾而忘反。\par
背夏浦而西思兮，哀故都之日远。\par
登大坟以远望兮，聊以舒吾忧心。\par
哀州土之平乐兮，悲江介之遗风。\par
当陵阳之焉至兮，淼南渡之焉如？\par
曾不知夏之为丘兮，孰两东门之可芜？\par
心不怡之长久兮，忧与愁其相接。\par
惟郢路之辽远兮，江与夏之不可涉。\par
忽若不信兮，至今九年而不复。\par
惨郁郁而不通兮，蹇侘傺而含慼。\par
外承欢之汋约兮，谌荏弱而难持。\par
忠湛湛而愿进兮，妒被离而鄣之。\par
尧舜之抗行兮，瞭杳杳而薄天。\par
众谗人之嫉妒兮，被以不慈之伪名。\par
憎愠惀之修美兮，好夫人之慷慨。\par
众踥蹀而日进兮，美超远而逾迈。\par
乱曰：\par
曼余目以流观兮，冀一反之何时？\par
鸟飞反故乡兮，狐死必首丘。\par
信非吾罪而弃逐兮，何日夜而忘之？\par

\section{抽思}
\textbf{九章　屈原}\par\vspace{0.5em}
心郁郁之忧思兮，独永叹乎增伤。\par
思蹇产之不释兮，曼遭夜之方长。\par
悲秋风之动容兮，何回极之浮浮。\par
数惟荪之多怒兮，伤余心之忧忧。\par
愿摇起而横奔兮，览民尤以自镇。\par
结微情以陈词兮，矫以遗夫美人。\par
昔君与我诚言兮，曰黄昏以为期。\par
羌中道而回畔兮，反既有此他志。\par
憍吾以其美好兮，览余以其修姱。\par
与余言而不信兮，盖为余而造怒。\par
愿承閒而自察兮，心震悼而不敢。\par
悲夷犹而冀进兮，心怛伤之憺憺。\par
兹历情以陈辞兮，荪详聋而不闻。\par
固切人之不媚兮，众果以我为患。\par
初吾所陈之耿著兮，岂至今其庸亡？\par
何独乐斯之謇謇兮？\par
愿荪美之可光。\par
望三王以为像兮，指彭咸以为仪。\par
夫何极而不至兮，故远闻而难亏。\par
善不由外来兮，名不可以虚作。\par
孰无施而有报兮，孰不实而有获？\par
少歌曰：与美人抽思兮，并日夜而无正。\par
憍吾以其美好兮，敖朕辞而不听。\par
倡曰：\par
有鸟自南兮，来集汉北。\par
好姱佳丽兮，牉独处此异域。\par
惸茕独而不群兮，又无良媒在其侧。\par
道卓远而日忘兮，愿自申而不得。\par
望北山而流涕兮，临流水而太息。\par
望孟夏之短夜兮，何晦明之若岁？\par
惟郢路之辽远兮，魂一夕而九逝。\par
曾不知路之曲直兮，南指月与列星。\par
愿径逝而未得兮，魂识路之营营。\par
何灵魂之信直兮，人之心不与吾心同！\par
理弱而媒不通兮，尚不知余之从容。\par
乱曰：\par
长濑湍流，溯江潭兮。\par
狂顾南行，聊以娱心兮。\par
轸石崴嵬，蹇吾愿兮。\par
超回志度，行隐进兮。\par
低徊夷犹，宿北姑兮。\par
烦冤瞀容，实沛徂兮。\par
愁叹苦神，灵遥思兮。\par
路远处幽，又无行媒兮。\par
道思作颂，聊以自救兮。\par
忧心不遂，斯言谁告兮。\par

\section{怀沙}
\textbf{九章　屈原}\par\vspace{0.5em}
滔滔孟夏兮，草木莽莽。\par
伤怀永哀兮，汩徂南土。\par
眴兮杳杳，孔静幽默。\par
郁结纡轸兮，离慜而长鞠。\par
抚情效志兮，冤屈而自抑。\par
刓方以为圜兮，常度未替。\par
易初本迪兮，君子所鄙。\par
章画志墨兮，前图未改。\par
内厚质正兮，大人所盛。\par
巧倕不斲兮，孰察其拨正。\par
玄文处幽兮，矇瞍谓之不章\par
离娄微睇兮，瞽以为无明。\par
变白以为黑兮，倒上以为下。\par
凤皇在笯兮，鸡鹜翔舞。\par
同糅玉石兮，一概而相量。\par
夫惟党人之鄙固兮，羌不知余之所臧。\par
任重载盛兮，陷滞而不济。\par
怀瑾握瑜兮，穷不知所示。\par
邑犬之群吠兮，吠所怪也。\par
非俊疑杰兮，固庸态也。\par
文质疏内兮，众不知余之异采。\par
材朴委积兮，莫知余之所有。\par
重仁袭义兮，谨厚以为丰。\par
重华不可遌兮，孰知余之从容！\par
古固有不并兮，岂知其何故也？\par
汤禹久远兮，邈而不可慕也？\par
惩违改忿兮，抑心而自强。\par
离慜而不迁兮，愿志之有像。\par
进路北次兮，日昧昧其将暮。\par
舒忧娱哀兮，限之以大故。\par
乱曰：\par
浩浩沅湘，分流汩兮。\par
修路幽蔽，道远忽兮。\par
怀质抱情，独无匹兮。\par
伯乐既没，骥焉程兮。\par
民生禀命，各有所错兮。\par
定心广志，余何所畏惧兮？\par
曾伤爰哀，永叹喟兮。\par
世浑浊莫吾知，人心不可谓兮。\par
知死不可让，愿勿爱兮。\par
明告君子，吾将以为类兮。\par

\section{思美人}
\textbf{九章　屈原}\par\vspace{0.5em}
思美人兮，揽涕而竚眙。\par
媒绝路阻兮，言不可结而诒。\par
蹇蹇之烦冤兮，陷滞而不发。\par
申旦以舒中情兮，志沉菀而莫达。\par
愿寄言于浮云兮，遇丰隆而不将。\par
因归鸟而致辞兮，羌迅高而难当。\par
高辛之灵盛兮，遭玄鸟而致诒。\par
欲变节以从俗兮，媿易初而屈志。\par
独历年而离愍兮，羌凭心犹未化。\par
宁隐闵而寿考兮，何变易之可为！\par
知前辙之不遂兮，未改此度。\par
车既覆而马颠兮，蹇独怀此异路。\par
勒骐骥而更驾兮，造父为我操之，\par
迁逡次而勿驱兮，聊假日以须是时。\par
指嶓冢之西隈兮，与纁黄以为期。\par
开春发岁兮，白日出之悠悠。\par
吾将荡志而愉乐兮，遵江夏以娱忧。\par
揽大薄之芳茝兮，搴长洲之宿莽。\par
惜吾不及古人兮，吾谁与玩此芳草？\par
解萹薄与杂菜兮，备以为交佩。\par
佩缤纷以缭转兮，遂萎绝而离异。\par
吾且儃徊以娱忧兮，观南人之变态。\par
窃快在中心兮，扬厥凭而不竢。\par
芳与泽其杂糅兮，羌芳华自中出。\par
纷郁郁其远蒸兮，满内而外扬。\par
情与质信可保兮，羌居蔽而闻章。\par
令薜荔以为理兮，惮举趾而缘木。\par
因芙蓉而为媒兮，惮褰裳而濡足。\par
登高吾不说兮，入下吾不能。\par
固朕形之不服兮，然容与而狐疑。\par
广遂前画兮，未改此度也。\par
命则处幽吾将罢兮，愿及白日之未暮也。\par
独茕茕而南行兮，思彭咸之故也。\par

\section{惜往日}
\textbf{九章　屈原}\par\vspace{0.5em}
惜往日之曾信兮，受命诏以昭时。\par
奉先功以照下兮，明法度之嫌疑。\par
国富强而法立兮，属贞臣而日竢。\par
秘密事之载心兮，虽过失犹弗治。\par
心纯庞而不泄兮，遭谗人而嫉之。\par
君含怒而待臣兮，不清澈其然否。\par
蔽晦君之聪明兮，虚惑误又以欺。\par
弗参验以考实兮，远迁臣而弗思。\par
信谗谀之浑浊兮，盛气志而过之。\par
何贞臣之无罪兮，被离谤而见尤。\par
惭光景之诚信兮，身幽隐而备之。\par
临沅湘之玄渊兮，遂自忍而沉流。\par
卒没身而绝名兮，惜壅君之不昭。\par
君无度而弗察兮，使芳草为薮幽。\par
焉舒情而抽信兮，恬死亡而不聊。\par
独障壅而弊隐兮，使贞臣为无由。\par
闻百里之为虏兮，伊尹烹于庖厨。\par
吕望屠于朝歌兮，宁戚歌而饭牛。\par
不逢汤武与桓缪兮，世孰云而知之。\par
吴信谗而弗味兮，子胥死而后忧。\par
介子忠而立枯兮，文君寤而追求。\par
封介山而为之禁兮，报大德之优游。\par
思久故之亲身兮，因缟素而哭之。\par
或忠信而死节兮，或訑谩而不疑。\par
弗省察而按实兮，听谗人之虚辞。\par
芳与泽其杂糅兮，孰申旦而别之？\par
何芳草之早殀兮，微霜降而下戒。\par
谅聪不明而蔽壅兮，使谗谀而日得。\par
自前世之嫉贤兮，谓蕙若其不可佩。\par
妒佳冶之芬芳兮，嫫母姣而自好。\par
虽有西施之美容兮，谗妒入以自代。\par
愿陈情以白行兮，得罪过之不意。\par
情冤见之日明兮，如列宿之错置。\par
乘骐骥而驰骋兮，无辔衔而自载\par
乘泛泭以下流兮，无舟楫而自备。\par
背法度而心治兮，辟与此其无异。\par
宁溘死而流亡兮，恐祸殃之有再。\par
不毕辞而赴渊兮，惜壅君之不识。\par

\section{橘颂}
\textbf{九章　屈原}\par\vspace{0.5em}
后皇嘉树，橘徕服兮。\par
受命不迁，生南国兮。\par
深固难徙，更壹志兮。\par
绿叶素荣，纷其可喜兮。\par
曾枝剡棘，圆果抟兮。\par
青黄杂糅，文章烂兮。\par
精色内白，类任道兮。\par
纷緼宜修，姱而不丑兮。\par
嗟尔幼志，有以异兮。\par
独立不迁，岂不可喜兮？\par
深固难徙，廓其无求兮。\par
苏世独立，横而不流兮。\par
闭心自慎，终不失过兮。\par
秉德无私，参天地兮。\par
愿岁并谢，与长友兮。\par
淑离不淫，梗其有理兮。\par
年岁虽少，可师长兮。\par
行比伯夷，置以为像兮。\par

\section{悲回风}
\textbf{九章　屈原}\par\vspace{0.5em}
悲回风之摇蕙兮，心冤结而内伤。\par
物有微而陨性兮，声有隐而先倡。\par
夫何彭咸之造思兮，暨志介而不忘！\par
万变其情岂可盖兮，孰虚伪之可长？\par
鸟兽鸣以号群兮，草苴比而不芳。\par
鱼葺鳞以自别兮，蛟龙隐其文章。\par
故荼荠不同亩兮，兰茝幽而独芳。\par
惟佳人之永都兮，更统世以自贶。\par
眇远志之所及兮，怜浮云之相羊。\par
介眇志之所惑兮，窃赋诗之所明。\par
惟佳人之独怀兮，折若椒以自处。\par
曾歔欷之嗟嗟兮，独隐伏而思虑。\par
涕泣交而凄凄兮，思不眠以至曙。\par
终长夜之曼曼兮，掩此哀而不去。\par
寤从容以周流兮，聊逍遥以自恃。\par
伤太息之愍怜兮，气于邑而不可止。\par
糺思心以为纕兮，编愁苦以为膺。\par
折若木以弊光兮，随飘风之所仍。\par
存彷佛而不见兮，心踊跃其若汤。\par
抚珮衽以案志兮，超惘惘而遂行。\par
岁曶曶其若颓兮，时亦冉冉而将至。\par
薠蘅槁而节离兮，芳以歇而不比。\par
怜思心之不可惩兮，证此言之不可聊。\par
宁溘死而流亡兮，不忍此心之常愁。\par
孤子吟而抆泪兮，放子出而不还。\par
孰能思而不隐兮，照彭咸之所闻。\par
登石峦以远望兮，路眇眇之默默。\par
入景响之无应兮，闻省想而不可得。\par
愁郁郁之无快兮，居戚戚而不可解。\par
心鞿羁而不开兮，气缭转而自缔。\par
穆眇眇之无垠兮，莽芒芒之无仪。\par
声有隐而相感兮，物有纯而不可为。\par
邈漫漫之不可量兮，缥绵绵之不可纡。\par
愁悄悄之常悲兮，翩冥冥之不可娱。\par
凌大波而流风兮，讬彭咸之所居。\par
上高岩之峭岸兮，处雌蜺之标颠。\par
据青冥而摅虹兮，遂儵忽而扪天。\par
吸湛露之浮源兮，漱凝霜之雰雰。\par
依风穴以自息兮，忽倾寤以婵媛。\par
冯昆仑以澂雾兮，隐渂山以清江。\par
惮涌湍之礚礚兮，听波声之汹汹。\par
纷容容之无经兮，罔芒芒之无纪。\par
轧洋洋之无从兮，驰委移之焉止？\par
漂翻翻其上下兮，翼遥遥其左右。\par
氾潏潏其前后兮，伴张驰之信期。\par
观炎气之相仍兮，窥烟液之所积。\par
悲霜雪之俱下兮，听潮水之相击。\par
借光景以往来兮，施黄棘之枉策。\par
求介子之所存兮，见伯夷之放迹。\par
心调度而弗去兮，刻著志之无适。\par
曰吾怨往昔之所冀兮，悼来者之悐悐。\par
浮江淮而入海兮，从子胥而自适。\par
望大河之洲渚兮，悲申徒之抗迹。\par
骤谏君而不听兮，重任石之何益？\par
心絓结而不解兮，思蹇产而不释。\par

\section{远游}
\textbf{远游　屈原}\par\vspace{0.5em}
悲时俗之迫阨兮，愿轻举而远游。\par
质菲薄而无因兮，焉讬乘而上浮？\par
遭沈浊而污秽兮，独郁结其谁语！\par
夜耿耿而不寐兮，魂营营而至曙。\par
惟天地之无穷兮，哀人生之长勤。\par
往者余弗及兮，来者吾不闻。\par
步徙倚而遥思兮，怊惝怳而乖怀。\par
意荒忽而流荡兮，心愁悽而增悲。\par
神倏忽而不反兮，形枯槁而独留。\par
内惟省以操端兮，求正气之所由。\par
漠虚静以恬愉兮，澹无为而自得。\par
闻赤松之清尘兮，愿承风乎遗则。\par
贵真人之休德兮，美往世之登仙\par
与化去而不见兮，名声著而日延。\par
奇傅说之讬辰星兮，羡韩众之得一。\par
形穆穆以浸远兮，离人群而遁逸。\par
因气变而遂曾举兮，忽神奔而鬼怪。\par
时仿佛以遥见兮，精晈晈以往来。\par
超氛埃而淑邮兮，终不反其故都。\par
免众患而不惧兮，世莫知其所如。\par
恐天时之代序兮，耀灵晔而西征。\par
微霜降而下沦兮，悼芳草之先蘦。\par
聊仿佯而逍遥兮，永历年而无成。\par
谁可与玩斯遗芳兮？\par
长向风而舒情。\par
高阳邈以远兮，余将焉所程？\par
重曰：\par
春秋忽其不淹兮，奚久留此故居。\par
轩辕不可攀援兮，吾将从王乔而娱戏。\par
餐六气而饮沆瀣兮，漱正阳而含朝霞。\par
保神明之清澄兮，精气入而麤秽除。\par
顺凯风以从游兮，至南巢而壹息。\par
见王子而宿之兮，审壹气之和德。\par
曰“道可受兮，不可传\par
其小无内兮，其大无垠。\par
毋滑而魂兮，彼将自然\par
壹气孔神兮，于中夜存。\par
虚以待之存，无为之先\par
庶类以成兮，此德之门。”\par
闻至贵而遂徂兮，忽乎吾将行。\par
仍羽人于丹丘，留不死之旧乡。\par
朝濯发于汤谷兮，夕晞余身兮九阳。\par
吸飞泉之微液兮，怀琬琰之华英。\par
玉色頩以脕颜兮，精醇粹而始壮。\par
质销铄以汋约兮，神要眇以淫放。\par
嘉南州之炎德兮，丽桂树之冬荣\par
山萧条而无兽兮，野寂漠其无人。\par
载营魄而登霞兮，掩浮云而上征。\par
命天阍其开关兮，排阊阖而望予。\par
召丰隆使先导兮，问太微之所居。\par
集重阳入帝宫兮，造旬始而观清都。\par
朝发轫于太仪兮，夕始临乎于微闾。\par
屯余车之万乘兮，纷容与而并驰。\par
驾八龙之婉婉兮，载云旗之逶蛇。\par
建雄虹之采旄兮，五色杂而炫耀。\par
服偃蹇以低昂兮，骖连蜷以骄骜。\par
骑胶葛以杂乱兮，斑漫衍而方行。\par
撰余辔而正策兮，吾将过乎句芒。\par
历太皓以右转兮，前飞廉以启路。\par
阳杲杲其未光兮，凌天地以径度。\par
风伯为余先驱兮，氛埃辟而清凉。\par
凤凰翼其承旂兮，遇蓐收乎西皇。\par
揽慧星以为旍兮，举斗柄以为麾。\par
叛陆离其上下兮，游惊雾之流波。\par
时暧曃其曭莽兮，召玄武而奔属。\par
后文昌使掌行兮，选署众神以并轂。\par
路漫漫其修远兮，徐弭节而高厉。\par
左雨师使径侍兮，右雷公以为卫。\par
欲度世以忘归兮，意姿睢以抯挢。\par
内欣欣而自美兮，聊媮娱以淫乐。\par
涉青云以汎滥游兮，忽临睨夫旧乡。\par
仆夫怀余心悲兮，边马顾而不行。\par
思旧故以想象兮，长太息而掩涕。\par
汜容与而遐举兮，聊抑志而自弭。\par
指炎神而直驰兮，吾将往乎南疑。\par
览方外之荒忽兮，沛罔瀁而自浮。\par
祝融戒而跸御兮，腾告鸾鸟迎宓妃。\par
张咸池奏承云兮，二女御九韶歌。\par
使湘灵鼓瑟兮，令海若舞冯夷。\par
玄螭虫象并出进兮，形蟉虯而逶蛇。\par
雌蜺便娟以增挠兮，鸾鸟轩翥而翔飞。\par
音乐博衍无终极兮，焉乃逝以徘徊。\par
舒并节以驰骛兮，逴绝垠乎寒门。\par
轶迅风于清源兮，从颛顼乎增冰。\par
历玄冥以邪径兮，乘间维以反顾。\par
召黔赢而见之兮，为余先乎平路。\par
经营四方兮，周流六漠。\par
上至列缺兮，降望大壑。\par
下峥嵘而无地兮，上寥廓而无天。\par
视倏忽而无见兮，听惝恍而无闻。\par
超无为以至清兮，与泰初而为邻。\par

\section{卜居}
\textbf{卜居　屈原}\par\vspace{0.5em}
屈原既放，三年不得复见。竭知尽忠而蔽障于谗。心烦虑乱，不知所从。乃往见太卜郑詹尹曰：“余有所疑，愿因先生决之。”詹尹乃端策拂龟，曰：“君将何以教之？”\par
屈原曰：“吾宁悃悃款款，朴以忠乎，将送往劳来，斯无穷乎？宁诛锄草茅以力耕乎，将游大人以成名乎？宁正言不讳以危身乎，将从俗富贵以偷生乎？宁超然高举以保真乎，将哫訾栗斯，喔咿儒儿，以事妇人乎？宁廉洁正直以自清乎，将突梯滑稽，如脂如韦，以洁楹乎？宁昂昂若千里之驹乎，将泛泛若水中之凫，与波上下，偷以全吾躯乎？宁与骐骥亢轭乎，将随驽马之迹乎？宁与黄鹄比翼乎，将与鸡鹜争食乎？此孰吉孰凶？何去何从？世溷浊而不清：蝉翼为重，千钧为轻黄钟毁弃，瓦釜雷鸣谗人高张，贤士无名。吁嗟默默兮，谁知吾之廉贞！”\par
詹尹乃释策而谢曰：“夫尺有所短，寸有所长物有所不足，智有所不明数有所不逮，神有所不通。用君之心，行君之意。龟策诚不能知此事。”\par

\section{渔父}
\textbf{渔父　屈原}\par\vspace{0.5em}
屈原既放，游于江潭，行吟泽畔，颜色憔悴，形容枯槁。渔父见而问之曰：“子非三闾大夫与？何故至于斯？”屈原曰：“举世皆浊我独清，众人皆醉我独醒，是以见放。”\par
渔父曰：“圣人不凝滞于物，而能与世推移。世人皆浊，何不淈其泥而扬其波？众人皆醉，何不餔其糟而歠其醨？何故深思高举，自令放为？”\par
屈原曰：“吾闻之，新沐者必弹冠，新浴者必振衣安能以身之察察，受物之汶汶者乎？宁赴湘流，葬于江鱼之腹中。安能以皓皓之白，而蒙世俗之尘埃乎?”\par
渔父莞尔而笑，鼓枻而去。\par
乃歌曰：“沧浪之水清兮，可以濯吾缨。沧浪之水浊兮，可以濯吾足。”遂去，不复与言。\par

\section{九辩}
\textbf{九辩　宋玉}\par\vspace{0.5em}
悲哉，秋之为气也！萧瑟兮草木摇落而变衰。\par
憭栗兮若在远行，登山临水兮送将归。\par
泬漻兮天高而气清，寂寥兮收潦而水清。\par
憯悽增欷兮，薄寒之中人，怆怳懭悢兮，去故而就新。\par
坎廪兮贫士失职而志不平，廓落兮羁旅而无友生，惆怅兮而私自怜！\par
燕翩翩其辞归兮，蝉寂漠而无声。\par
雁廱廱而南游兮，鹍鸡啁哳而悲鸣。\par
独申旦而不寐兮，哀蟋蟀之宵征。\par
时亹亹而过中兮，蹇淹留而无成。\par
悲忧穷戚兮独处廓，有美一人兮心不绎。\par
去乡离家兮来远客，超逍遥兮今焉薄！\par
专思君兮不可化，君不知兮可奈何！\par
蓄怨兮积思，心烦憺兮忘食事。\par
原一见兮道余意，君之心兮与余异。\par
车既驾兮朅而归，不得见兮心伤悲。\par
倚结軨兮长太息，涕潺湲兮下霑轼。\par
忼慨绝兮不得，中瞀乱兮迷惑。\par
私自怜兮何极？\par
心怦怦兮谅直。\par
皇天平分四时兮，窃独悲此凛秋。\par
白露既下百草兮，奄离披此梧楸。\par
去白日之昭昭兮，袭长夜之悠悠。\par
离芳蔼之方壮兮，余萎约而悲愁。\par
秋既先戒以白露兮，冬又申之以严霜。\par
收恢台之孟夏兮，然欿傺而沉藏。\par
叶菸邑而无色兮，枝烦挐而交横。\par
颜淫溢而将罢兮，柯仿佛而萎黄。\par
萷櫹椮之可哀兮，形销铄而瘀伤。\par
惟其纷糅而将落兮，恨其失时而无当。\par
揽騑辔而下节兮，聊逍遥以相佯。\par
岁忽忽而遒尽兮，恐余寿之弗将。\par
悼余生之不时兮，逢此世之俇攘。\par
澹容与而独倚兮，蟋蟀鸣此西堂。\par
心怵惕而震荡兮，何所忧之多方。\par
卬明月而太息兮，步列星而极明。\par
窃悲夫蕙华之曾敷兮，纷旖旎乎都房。\par
何曾华之无实兮，从风雨而飞飏！\par
以为君独服此蕙兮，羌无以异于众芳。\par
闵奇思之不通兮，将去君而高翔。\par
心闵怜之惨悽兮，愿一见而有明。\par
重无怨而生离兮，中结轸而增伤。\par
岂不郁陶而思君兮？\par
君之门以九重！猛犬狺狺而迎吠兮，关梁闭而不通。\par
皇天淫溢而秋霖兮，后土何时而得漧？\par
塊独守此无泽兮，仰浮云而永叹！何时俗之工巧兮？\par
背绳墨而改错！郤骐骥而不乘兮，策驽骀而取路。\par
当世岂无骐骥兮，诚莫之能善御。\par
见执辔者非其人兮，故駶跳而远去。\par
凫雁皆唼夫梁藻兮，凤愈飘翔而高举。\par
圜凿而方枘兮，吾固知其鉏铻而难入。\par
众鸟皆有所登棲兮，凤独遑遑而无所集。\par
原衔枚而无言兮，尝被君之渥洽。\par
太公九十乃显荣兮，诚未遇其匹合。\par
谓骐骥兮安归？\par
谓凤皇兮安棲？\par
变古易俗兮世衰，今之相者兮举肥。\par
骐骥伏匿而不见兮，凤皇高飞而不下。\par
鸟兽犹知怀德兮，何云贤士之不处？\par
骥不骤进而求服兮，凤亦不贪餧而妄食。\par
君弃远而不察兮，虽原忠其焉得？\par
欲寂漠而绝端兮，窃不敢忘初之厚德。\par
独悲愁其伤人兮，冯郁郁其何极？\par
霜露惨悽而交下兮，心尚幸其弗济。\par
霰雪雰糅其增加兮，乃知遭命之将至。\par
原徼幸而有待兮，泊莽莽与野草同死。\par
原自往而径游兮，路壅绝而不通。\par
欲循道而平驱兮，又未知其所从。\par
然中路而迷惑兮，自压桉而学诵。\par
性愚陋以褊浅兮，信未达乎从容。\par
窃美申包胥之气盛兮，恐时世之不固。\par
何时俗之工巧兮？\par
灭规矩而改凿！独耿介而不随兮，原慕先圣之遗教。\par
处浊世而显荣兮，非余心之所乐。\par
与其无义而有名兮，宁穷处而守高。\par
食不媮而为饱兮，衣不苟而为温。\par
窃慕诗人之遗风兮，原讬志乎素餐。\par
蹇充倔而无端兮，泊莽莽而无垠。\par
无衣裘以御冬兮，恐溘死不得见乎阳春。\par
靓杪秋之遥夜兮，心缭悷而有哀。\par
春秋逴逴而日高兮，然惆怅而自悲。\par
四时递来而卒岁兮，阴阳不可与俪偕。\par
白日晼晚其将入兮，明月销铄而减毁。\par
岁忽忽而遒尽兮，老冉冉而愈弛。\par
心摇悦而日幸兮，然怊怅而无冀。\par
中憯恻之悽怆兮，长太息而增欷。\par
年洋洋以日往兮，老嵺廓而无处。\par
事亹亹而觊进兮，蹇淹留而踌躇。\par
何氾滥之浮云兮？\par
猋壅蔽此明月。\par
忠昭昭而原见兮，然霠曀而莫达。\par
原皓日之显行兮，云蒙蒙而蔽之。\par
窃不自聊而原忠兮，或黕点而汙之。\par
尧舜之抗行兮，瞭冥冥而薄天。\par
何险巇之嫉妒兮？\par
被以不慈之伪名。\par
彼日月之照明兮，尚黯黮而有瑕。\par
何况一国之事兮，亦多端而胶加。\par
被荷裯之晏晏兮，然潢洋而不可带。\par
既骄美而伐武兮，负左右之耿介。\par
憎愠惀之修美兮，好夫人之慷慨。\par
众踥蹀而日进兮，美超远而逾迈。\par
农夫辍耕而容与兮，恐田野之芜秽。\par
事緜緜而多私兮，窃悼後之危败。\par
世雷同而炫曜兮，何毁誉之昧昧！\par
今修饰而窥镜兮，後尚可以竄藏。\par
愿寄言夫流星兮，羌倏忽而难当。\par
卒壅蔽此浮云，下暗漠而无光。\par
尧舜皆有所举任兮，故高枕而自适。\par
谅无怨于天下兮，心焉取此怵惕？\par
乘骐骥之浏浏兮，驭安用夫强策？\par
谅城郭之不足恃兮，虽重介之何益？\par
邅翼翼而无终兮，忳惛惛而愁约。\par
生天地之若过兮，功不成而无嶜。\par
原沉滞而不见兮，尚欲布名乎天下。\par
然潢洋而不遇兮，直怐愗而自苦。\par
莽洋洋而无极兮，忽翱翔之焉薄？\par
国有骥而不知乘兮，焉皇皇而更索？\par
宁戚讴于车下兮，桓公闻而知之。\par
无伯乐之相善兮，今谁使乎誉之？\par
罔流涕以聊虑兮，惟著意而得之。\par
纷纯纯之愿忠兮，妒被离而鄣之。\par
原赐不肖之躯而别离兮，放游志乎云中。\par
乘精气之抟抟兮，骛诸神之湛湛。\par
骖白霓之習習兮，历群灵之丰丰。\par
左硃雀之茇茇兮，右苍龙之躣躣。\par
属雷师之阗阗兮，通飞廉之衙衙。\par
前轻辌之锵锵兮，后辎乘之从从。\par
载云旗之委蛇兮，扈屯骑之容容。\par
计专专之不可化兮，原遂推而为臧。\par
赖皇天之厚德兮，还及君之无恙！\par

\section{招魂}
\textbf{招魂　宋玉}\par\vspace{0.5em}
朕幼清以廉洁兮，身服义而未沫。\par
主此盛德兮，牵于俗而芜秽。\par
上无所考此盛德兮，长离殃而愁苦。\par
帝告巫阳曰：“有人在下，我欲辅之。魂魄离散，汝筮予之。”\par
巫阳对曰：“掌梦！上帝其难从。若必筮予之，恐后之谢，不能复用。”\par
巫阳焉乃下招曰：魂兮归来！去君之恒干，何为四方些？\par
舍君之乐处，而离彼不祥些！魂兮归来！东方不可以讬些。\par
长人千仞，惟魂是索些。\par
十日代出，流金铄石些。\par
彼皆习之，魂往必释些。\par
归来兮！不可以讬些。\par
魂兮归来！南方不可以止些。\par
雕题黑齿，得人肉以祀，以其骨为醢些。\par
蝮蛇蓁蓁，封狐千里些。\par
雄虺九首，往来倏忽，吞人以益其心些。\par
归来兮！不可久淫些。\par
魂兮归来！西方之害，流沙千里些。\par
旋入雷渊，爢散而不可止些。\par
幸而得脱，其外旷宇些。\par
赤蚁若象，玄蜂若壶些。\par
五谷不生，丛菅是食些。\par
其土烂人，求水无所得些。\par
彷徉无所倚，广大无所极些。\par
归来兮！恐自遗贼些。\par
魂兮归来！北方不可以止些。\par
增冰峨峨，飞雪千里些。\par
归来兮！不可以久些。\par
魂兮归来！君无上天些。\par
虎豹九关，啄害下人些。\par
一夫九首，拔木九千些。\par
豺狼从目，往来侁侁些。\par
悬人以嬉，投之深渊些。\par
致命于帝，然后得瞑些。\par
归来！往恐危身些。\par
魂兮归来！君无下此幽都些。\par
土伯九约，其角觺觺些。\par
敦脄血拇，逐人伂駓駓些。\par
参目虎首，其身若牛些。\par
此皆甘人，归来！恐自遗灾些。\par
魂兮归来！入修门些。\par
工祝招君，背行先些。\par
秦篝齐缕，郑绵络些。\par
招具该备，永啸呼些。\par
魂兮归来！反故居些。\par
天地四方，多贼奸些。\par
像设君室，静闲安些。\par
高堂邃宇，槛层轩些。\par
层台累榭，临高山些。\par
网户朱缀，刻方连些。\par
冬有穾厦，夏室寒些。\par
川谷径复，流潺湲些。\par
光风转蕙，氾崇兰些。\par
经堂入奥，朱尘筵些。\par
砥室翠翘，挂曲琼些。\par
翡翠珠被，烂齐光些。\par
蒻阿拂壁，罗帱张些。\par
纂组绮缟，结琦璜些。\par
室中之观，多珍怪些。\par
兰膏明烛，华容备些。\par
二八侍宿，射递代些。\par
九侯淑女，多迅众些。\par
盛鬋不同制，实满宫些。\par
容态好比，顺弥代些。\par
弱颜固植，謇其有意些。\par
姱容修态，絚洞房些。\par
蛾眉曼睩，目腾光些。\par
靡颜腻理，遗视矊些。\par
离榭修幕，侍君之闲些。\par
悲帷翠帐，饰高堂些。\par
红壁沙版，玄玉梁些。\par
仰观刻桷，画龙蛇些。\par
坐堂伏槛，临曲池些。\par
芙蓉始发，杂芰荷些。\par
紫茎屏风，文缘波些。\par
文异豹饰，侍陂陁些。\par
轩辌既低，步骑罗些。\par
兰薄户树，琼木篱些。\par
魂兮归来！何远为些？\par
室家遂宗，食多方些。\par
稻粢穱麦，挐黄梁些。\par
大苦醎酸，辛甘行些。\par
肥牛之腱，臑若芳些。\par
和酸若苦，陈吴羹些。\par
胹鳖炮羔，有柘浆些。\par
鹄酸臇凫，煎鸿鸧些。\par
露鸡臛蠵，厉而不爽些。\par
粔籹蜜饵，有餦餭些。\par
瑶浆蜜勺，实羽觞些。\par
挫糟冻饮，酎清凉些。\par
华酌既陈，有琼浆些。\par
归来反故室，敬而无妨些。\par
肴羞未通，女乐罗些。\par
敶钟按鼓，造新歌些。\par
《涉江》《采菱》，发《扬荷》些。\par
美人既醉，朱颜酡些。\par
嬉光眇视，目曾波些。\par
被文服纤，丽而不奇些。\par
长发曼鬋，艳陆离些。\par
二八齐容，起郑舞些。\par
衽若交竿，抚案下些。\par
竽瑟狂会，搷鸣鼓些。\par
宫庭震惊，发<激楚>些。\par
吴歈蔡讴，奏大吕些。\par
士女杂坐，乱而不分些。\par
放敶组缨，班其相纷些。\par
郑卫妖玩，来杂陈些。\par
《激楚》之结，独秀先些。\par
菎蔽象棋，有六簙些。\par
分曹并进，遒相迫些。\par
成枭而牟，呼五白些。\par
晋制犀比，费白日些。\par
铿钟摇簴，揳梓瑟些。\par
娱酒不废，沈日夜些。\par
兰膏明烛，华灯错些。\par
结撰至思，兰芳假些。\par
人有所极，同心赋些。\par
酎饮尽欢，乐先故些。\par
魂兮归来！反故居些。\par
乱曰：献岁发春兮，汨吾南征。\par
菉蘋齐叶兮，白芷生。\par
路贯庐江兮，左长薄。\par
倚沼畦瀛兮，遥望博。\par
青骊结驷兮，齐千乘。\par
悬火延起兮，玄颜烝。\par
步及骤处兮，诱骋先。\par
抑骛若通兮，引车右还。\par
与王趋梦兮，课后先。\par
君王亲发兮，惮青兕。\par
朱明承夜兮，时不可以淹。\par
皋兰被径兮，斯路渐。\par
湛湛江水兮，上有枫。\par
目极千里兮，伤春心。\par
魂兮归来，哀江南！\par

\section{大招}
\textbf{大招　景差}\par\vspace{0.5em}
青春受谢，白日昭只。\par
春气奋发，万物遽只。\par
冥凌浃行，魂无逃只。\par
魂魄归来！无远遥只。\par
魂乎归来！无东无西，无南无北只。\par
东有大海，溺水浟浟只。\par
螭龙并流，上下悠悠只。\par
雾雨淫淫，白皓胶只。\par
魂乎无东！汤谷寂寥只。\par
魂乎无南！南有炎火千里，蝮蛇蜒只。\par
山林险隘，虎豹蜿只。\par
鰅鳙短狐，王虺骞只。\par
魂乎无南！蜮伤躬只\par
魂乎无西！西方流沙，漭洋洋只。\par
豕首纵目，被发鬤只。\par
长爪踞牙，诶笑狂只。\par
魂乎无西！多害伤只。\par
魂乎无北！北有寒山，趠龙赩只。\par
代水不可涉，深不可测只。\par
天白颢颢，寒凝凝只。\par
魂乎无往！盈北极只。\par
魂魄归来！闲以静只。\par
自恣荆楚，安以定只。\par
逞志究欲，心意安只。\par
穷身永乐，年寿延只。\par
魂乎归来！乐不可言只。\par
五谷六仞，设菰梁只。\par
鼎臑盈望，和致芳只。\par
内鸧鸽鹄，味豺羹只。\par
魂乎归来！恣所尝只。\par
鲜蠵甘鸡，和楚酪只。\par
醢豚苦狗，脍苴蒪只。\par
吴酸蒿蒌，不沾薄只。\par
魂兮归来！恣所择只。\par
炙鸹烝凫，煔鹑敶只。\par
煎鰿膗雀，遽爽存只。\par
魂乎归来！丽以先只。\par
四酎并孰，不涩嗌只。\par
清馨冻饮，不歠役只。\par
吴醴白蘖，和楚沥只。\par
魂乎归来！不遽惕只。\par
代秦郑卫，鸣竽张只。\par
伏戏《驾辩》，楚《劳商》只。\par
讴和《扬阿》，赵萧倡只。\par
魂乎归来！定空桑只。\par
二八接舞，投诗赋只。\par
叩钟调磬，娱人乱只。\par
四上竞气，极声变只。\par
魂乎归来！听歌譔只。\par
朱唇皓齿，嫭以姱只。\par
比德好闲，习以都只。\par
丰肉微骨，调以娱只。\par
魂乎归来！安以舒只。\par
嫮目宜笑，娥眉曼只。\par
容则秀雅，稚朱颜只。\par
魂乎归来！静以安只。\par
姱修滂浩，丽以佳只。\par
曾颊倚耳，曲眉规只。\par
滂心绰态，姣丽施只。\par
小腰秀颈，若鲜卑只。\par
魂乎归来！思怨移只。\par
易中利心，以动作只。\par
粉白黛黑，施芳泽只。\par
长袂拂面，善留客只。\par
魂乎归来！以娱昔只。\par
青色直眉，美目媔只。\par
靥辅奇牙，宜笑嘕只。\par
丰肉微骨，体便娟只。\par
魂乎归来！恣所便只。\par
夏屋广大，沙堂秀只。\par
南房小坛，观绝霤只。\par
曲屋步壛，宜扰畜只。\par
腾驾步游，猎春囿只。\par
琼轂错衡，英华假只。\par
茝兰桂树，郁弥路只。\par
魂乎归来！恣志虑只。\par
孔雀盈园，畜鸾皇只！\par
鵾鸿群晨，杂鶖鸧只。\par
鸿鹄代游，曼骕驦只。\par
魂乎归来！凤凰翔只。\par
曼泽怡面，血气盛只。\par
永宜厥身，保寿命只。\par
室家盈廷，爵禄盛只。\par
魂乎归来！居室定只。\par
接径千里，出若云只。\par
三圭重侯，听类神只。\par
察笃夭隐，孤寡存只。\par
魂兮归来！正始昆只。\par
田邑千畛，人阜昌只。\par
美冒众流，德泽章只。\par
先威后文，善美明只。\par
魂乎归来！赏罚当只。\par
名声若日，照四海只。\par
德誉配天，万民理只。\par
北至幽陵，南交阯只。\par
西薄羊肠，东穷海只。\par
魂乎归来！尚贤士只。\par
发政献行，禁苛暴只。\par
举杰压陛，诛讥罢只。\par
直赢在位，近禹麾只。\par
豪杰执政，流泽施只。\par
魂乎来归！国家为只。\par
雄雄赫赫，天德明只。\par
三公穆穆，登降堂只。\par
诸侯毕极，立九卿只。\par
昭质既设，大侯张只。\par
执弓挟矢，揖辞让只。\par
魂乎来归！尚三王只。\par

\section{惜誓}
\textbf{惜誓　贾谊}\par\vspace{0.5em}
惜余年老而日衰兮，岁忽忽而不反。\par
登苍天而高举兮，历众山而日远。\par
观江河之纡曲兮，离四海之霑濡。\par
攀北极而一息兮，吸沆瀣以充虚。\par
飞朱鸟使先驱兮，驾太一之象舆。\par
苍龙蚴虯于左骖兮，白虎骋而为右騑。\par
建日月以为盖兮，载玉女于後车。\par
驰骛于杳冥之中兮，休息虖昆仑之墟。\par
乐穷极而不厌兮，愿从容虖神明。\par
涉丹水而驰骋兮，右大夏之遗风。\par
黄鹄之一举兮，知山川之纡曲。\par
再举兮，睹天地之圜方。\par
临中国之众人兮，讬回飙乎尚羊。\par
乃至少原之野兮，赤松、王乔皆在旁。\par
二子拥瑟而调均兮，余因称乎清商。\par
澹然而自乐兮，吸众气而翱翔。\par
念我长生而久仙兮，不如反余之故乡。\par
黄鹄後时而寄处兮，鸱枭群而制之。\par
神龙失水而陆居兮，为蝼蚁之所裁。\par
夫黄鹄神龙犹如此兮，况贤者之逢乱世哉。\par
寿冉冉而日衰兮，固儃回而不息。\par
俗流从而不止兮，众枉聚而矫直。\par
或偷合而苟进兮，或隐居而深藏。\par
苦称量之不审兮，同权概而就衡。\par
或推迻而苟容兮，或直言之谔謣。\par
伤诚是之不察兮，并纫茅丝以为索。\par
方世俗之幽昏兮，眩白黑之美恶。\par
放山渊之龟玉兮，相与贵夫砾石。\par
梅伯数谏而至醢兮，来革顺志而用国。\par
悲仁人之尽节兮，反为小人之所贼。\par
比干忠谏而剖心兮，箕子被发而佯狂。\par
水背流而源竭兮，木去根而不长。\par
非重躯以虑难兮，惜伤身之无功。\par
已矣哉！独不见夫鸾凤之高翔兮，乃集大皇之野。\par
循四极而回周兮，见盛德而後下。\par
彼圣人之神德兮，远浊世而自藏。\par
使麒麟可得羁而係兮，又何以异虖犬羊？\par

\section{招隐士}
\textbf{招隐士　淮南小山}\par\vspace{0.5em}
桂树丛生兮山之幽，偃蹇连蜷兮枝相缭。\par
山气巄嵷兮石嵯峨，溪谷崭岩兮水曾波。\par
猿狖群啸兮虎豹嗥，攀援桂枝兮聊淹留。\par
王孙游兮不归，春草生兮萋萋。\par
岁暮兮不自聊，蟪蛄鸣兮啾啾。\par
坱兮轧，山曲岪，心淹留兮恫慌忽。\par
罔兮沕，憭兮栗，虎豹穴。\par
丛薄深林兮，人上栗。\par
嵚岑碕礒兮，碅磳磈硊\par
树轮相纠兮，林木茷骫。\par
青莎杂树兮，薠草靃靡\par
白鹿麏麚兮，或腾或倚。\par
状貌崟崟兮峨峨，凄凄兮漇漇。\par
猕猴兮熊罴，慕类兮以悲\par
攀援桂枝兮聊淹留。\par
虎豹斗兮熊罴咆，禽兽骇兮亡其曹。\par
王孙兮归来，山中兮不可以久留。\par

\section{初放}
\textbf{七谏　东方朔}\par\vspace{0.5em}
平生于国兮，长于原野。\par
言语讷譅兮，又无彊辅。\par
浅智褊能兮，闻见又寡。\par
数言便事兮，见怨门下。\par
王不察其长利兮，卒见弃乎原野。\par
伏念思过兮，无可改者。\par
群众成朋兮，上浸以惑。\par
巧佞在前兮，贤者灭息。\par
尧、舜圣已没兮，孰为忠直？\par
高山崔巍兮，水流汤汤。\par
死日将至兮，与麋鹿同坑。\par
塊兮鞠，当道宿，举世皆然兮，余将谁告？\par
斥逐鸿鹄兮，近习鸱枭，斩伐橘柚兮，列树苦桃。\par
便娟之修竹兮，寄生乎江潭。\par
上葳蕤而防露兮，下泠泠而来风。\par
孰知其不合兮，若竹柏之异心。\par
往者不可及兮，来者不可待。\par
悠悠苍天兮，莫我振理。\par
窃怨君之不寤兮，吾独死而後已。\par

\section{沉江}
\textbf{七谏　东方朔}\par\vspace{0.5em}
惟往古之得失兮，览私微之所伤。\par
尧舜圣而慈仁兮，後世称而弗忘。\par
齐桓失于专任兮，夷吾忠而名彰。\par
晋献惑于孋姬兮，申生孝而被殃。\par
偃王行其仁义兮，荆文寤而徐亡。\par
纣暴虐以失位兮，周得佐乎吕望。\par
修往古以行恩兮，封比干之丘垄。\par
贤俊慕而自附兮，日浸淫而合同。\par
明法令而修理兮，兰芷幽而有芳。\par
苦众人之妒予兮，箕子寤而佯狂。\par
不顾地以贪名兮，心怫郁而内伤。\par
联蕙芷以为佩兮，过鲍肆而失香。\par
正臣端其操行兮，反离谤而见攘。\par
世俗更而变化兮，伯夷饿于首阳。\par
独廉洁而不容兮，叔齐久而逾明。\par
浮云陈而蔽晦兮，使日月乎无光。\par
忠臣贞而欲谏兮，谗谀毁而在旁。\par
秋草荣其将实兮，微霜下而夜降。\par
商风肃而害生兮，百草育而不长。\par
众并谐以妒贤兮，孤圣特而易伤。\par
怀计谋而不见用兮，岩穴处而隐藏。\par
成功隳而不卒兮，子胥死而不葬。\par
世从俗而变化兮，随风靡而成行。\par
信直退而毁败兮，虚伪进而得当。\par
追悔过之无及兮，岂尽忠而有功。\par
废制度而不用兮，务行私而去公。\par
终不变而死节兮，惜年齿之未央。\par
将方舟而下流兮，冀幸君之发矇。\par
痛忠言之逆耳兮，恨申子之沉江。\par
愿悉心之所闻兮，遭值君之不聪。\par
不开寤而难道兮，不别横之与纵。\par
听奸臣之浮说兮，绝国家之久长。\par
灭规矩而不用兮，背绳墨之正方。\par
离忧患而乃寤兮，若纵火于秋蓬。\par
业失之而不救兮，尚何论乎祸凶。\par
彼离畔而朋党兮，独行之士其何望？\par
日渐染而不自知兮，秋毫微哉而变容。\par
众轻积而折轴兮，原咎杂而累重。\par
赴湘沅之流澌兮，恐逐波而复东。\par
怀沙砾而自沉兮，不忍见君之蔽壅。\par

\section{怨世}
\textbf{七谏　东方朔}\par\vspace{0.5em}
世沉淖而难论兮，俗岒峨而嵾嵯。\par
清泠泠而歼灭兮，溷湛湛而日多。\par
枭鸮既以成群兮，玄鹤弭翼而屏移。\par
蓬艾亲入御于床笫兮，马兰踸踔而日加。\par
弃捐药芷与杜衡兮，余柰世之不知芳何？\par
何周道之平易兮，然芜秽而险戏。\par
高阳无故而委尘兮，唐虞点灼而毁议。\par
谁使正其真是兮，虽有八师而不可为。\par
皇天保其高兮，后土持其久。\par
服清白以逍遥兮，偏与乎玄英异色。\par
西施媞媞而不得见兮，嫫母勃屑而日侍。\par
桂蠹不知所淹留兮，蓼虫不知徙乎葵菜。\par
处湣湣之浊世兮，今安所达乎吾志。\par
意有所载而远逝兮，固非众人之所识。\par
骥踌躇于弊輂兮，遇孙阳而得代。\par
吕望穷困而不聊生兮，遭周文而舒志。\par
宁戚饭牛而商歌兮，桓公闻而弗置。\par
路室女之方桑兮，孔子过之以自侍。\par
吾独乖剌而无当兮，心悼怵而耄思。\par
思比干之恲恲兮，哀子胥之慎事。\par
悲楚人之和氏兮，献宝玉以为石。\par
遇厉武之不察兮，羌两足以毕斮。\par
小人之居势兮，视忠正之何若？\par
改前圣之法度兮，喜嗫嚅而妄作。\par
亲谗谀而疏贤圣兮，讼谓闾娵为丑恶。\par
愉近习而蔽远兮，孰知察其黑白？\par
卒不得效其心容兮，安眇眇而无所归薄。\par
专精爽以自明兮，晦冥冥而壅蔽。\par
年既已过太半兮，然埳轲而留滞。\par
欲高飞而远集兮，恐离罔而灭败。\par
独冤抑而无极兮，伤精神而寿夭。\par
皇天既不纯命兮，余生终无所依。\par
愿自沉于江流兮，绝横流而径逝。\par
宁为江海之泥涂兮，安能久见此浊世？\par

\section{怨思}
\textbf{七谏　东方朔}\par\vspace{0.5em}
贤士穷而隐处兮，廉方正而不容。\par
子胥谏而靡躯兮，比干忠而剖心。\par
子推自割而飤君兮，德日忘而怨深。\par
行明白而曰黑兮，荆棘聚而成林。\par
江离弃于穷巷兮，蒺藜蔓乎东厢。\par
贤者蔽而不见兮，谗谀进而相朋。\par
枭鸮并进而俱鸣兮，凤皇飞而高翔。\par
原壹往而径逝兮，道壅绝而不通。\par

\section{自悲}
\textbf{七谏　东方朔}\par\vspace{0.5em}
居愁懃其谁告兮，独永思而忧悲。\par
内自省而不惭兮，操愈坚而不衰。\par
隐三年而无决兮，岁忽忽其若颓。\par
怜余身不足以卒意兮，冀一见而复归。\par
哀人事之不幸兮，属天命而委之咸池。\par
身被疾而不闲兮，心沸热其若汤。\par
冰炭不可以相并兮，吾固知乎命之不长。\par
哀独苦死之无乐兮，惜予年之未央。\par
悲不反余之所居兮，恨离予之故乡。\par
鸟兽惊而失群兮，犹高飞而哀鸣。\par
狐死必首丘兮，夫人孰能不反其真情？\par
故人疏而日忘兮，新人近而俞好。\par
莫能行于杳冥兮，孰能施于无报？\par
苦众人之皆然兮，乘回风而远游。\par
凌恆山其若陋兮，聊愉娱以忘忧。\par
悲虚言之无实兮，苦众口之铄金。\par
过故乡而一顾兮，泣歔欷而霑衿。\par
厌白玉以为面兮，怀琬琰以为心。\par
邪气入而感内兮，施玉色而外淫。\par
何青云之流澜兮，微霜降之蒙蒙。\par
徐风至而徘徊兮，疾风过之汤汤。\par
闻南籓乐而欲往兮，至会稽而且止。\par
见韩众而宿之兮，问天道之所在？\par
借浮云以送予兮，载雌霓而为旌。\par
驾青龙以驰骛兮，班衍衍之冥冥。\par
忽容容其安之兮，超慌忽其焉如？\par
苦众人之难信兮，愿离群而远举。\par
登峦山而远望兮，好桂树之冬荣。\par
观天火之炎炀兮，听大壑之波声。\par
引八维以自道兮，含沆瀣以长生。\par
居不乐以时思兮，食草木之秋实。\par
饮菌若之朝露兮，构桂木而为室。\par
杂橘柚以为囿兮，列新夷与椒桢。\par
鹍鹤孤而夜号兮，哀居者之诚贞。\par

\section{哀命}
\textbf{七谏　东方朔}\par\vspace{0.5em}
哀时命之不合兮，伤楚国之多忧。\par
内怀情之洁白兮，遭乱世而离尤。\par
恶耿介之直行兮，世溷浊而不知。\par
何君臣之相失兮，上沅湘而分离。\par
测汨罗之湘水兮，知时固而不反。\par
伤离散之交乱兮，遂侧身而既远。\par
处玄舍之幽门兮，穴岩石而窟伏。\par
从水蛟而为徙兮，与神龙乎休息。\par
何山石之崭岩兮，灵魂屈而偃蹇。\par
含素水而蒙深兮，日眇眇而既远。\par
哀形体之离解兮，神罔两而无舍。\par
惟椒兰之不反兮，魂迷惑而不知路。\par
愿无过之设行兮，虽灭没之自乐。\par
痛楚国之流亡兮，哀灵修之过到。\par
固时俗之溷浊兮，志瞀迷而不知路。\par
念私门之正匠兮，遥涉江而远去。\par
念女嬃之婵媛兮，涕泣流乎于悒。\par
我决死而不生兮，虽重追吾何及。\par
戏疾濑之素水兮，望高山之蹇产。\par
哀高丘之赤岸兮，遂没身而不反。\par

\section{谬谏}
\textbf{七谏　东方朔}\par\vspace{0.5em}
怨灵修之浩荡兮，夫何执操之不固？\par
悲太山之为隍兮，孰江河之可涸？\par
愿承闲而效志兮，恐犯忌而干讳。\par
卒抚情以寂寞兮，然怊怅而自悲。\par
玉与石其同匮兮，贯鱼眼与珠玑。\par
驽骏杂而不分兮，服罢牛而骖骥。\par
年滔滔而自远兮，寿冉冉而愈衰。\par
心悇憛而烦冤兮，蹇超摇而无冀。\par
固时俗之工巧兮，灭规矩而改错。\par
郤骐骥而不乘兮，策驽骀而取路。\par
当世岂无骐骥兮，诚无王良之善驭。\par
见执辔者非其人兮，故驹跳而远去。\par
不量凿而正枘兮，恐矩矱之不同。\par
不论世而高举兮，恐操行之不调。\par
弧弓弛而不张兮，孰云知其所至？\par
无倾危之患难兮，焉知贤士之所死？\par
俗推佞而进富兮，节行张而不著。\par
贤良蔽而不群兮，朋曹比而党誉。\par
邪说饰而多曲兮，正法弧而不公。\par
直士隐而避匿兮，谗谀登乎明堂。\par
弃彭咸之娱乐兮，灭巧倕之绳墨。\par
菎蕗杂于黀蒸兮，机蓬矢以射革。\par
驾蹇驴而无策兮，又何路之能极？\par
以直鍼而为钓兮，又何鱼之能得？\par
伯牙之绝弦兮，无锺子期而听之。\par
和抱璞而泣血兮，安得良工而剖之？\par
同音者相和兮，同类者相似。\par
飞鸟号其群兮，鹿鸣求其友。\par
故叩宫而宫应兮，弹角而角动。\par
虎啸而谷风至兮，龙举而景云往。\par
音声之相和兮，言物类之相感也。\par
夫方圜之异形兮，势不可以相错。\par
列子隐身而穷处兮，世莫可以寄讬。\par
众鸟皆有行列兮，凤独翔翔而无所薄。\par
经浊世而不得志兮，愿侧身岩穴而自讬。\par
欲阖口而无言兮，尝被君之厚德。\par
独便悁而怀毒兮，愁郁郁之焉极？\par
念三年之积思兮，愿壹见而陈辞。\par
不及君而骋说兮，世孰可为明之？\par
身寝疾而日愁兮，情沉抑而不扬。\par
众人莫可与论道兮，悲精神之不通。\par
乱曰：鸾皇孔凤日以远兮，畜凫驾鹅。\par
鸡鹜满堂坛兮，鼉黽游乎华池。\par
要褭奔亡兮，腾驾橐驼。\par
铅刀进御兮，遥弃太阿。\par
拔搴玄芝兮，列树芋荷。\par
橘柚萎枯兮，苦李旖旎。\par
甂瓯登于明堂兮，周鼎潜潜乎深渊。\par
自古而固然兮，吾又何怨乎今之人。\par

\section{哀时命}
\textbf{哀时命　庄忌}\par\vspace{0.5em}
哀时命之不及古人兮，夫何予生之不遘时！\par
往者不可扳援兮，徠者不可与期。\par
志憾恨而不逞兮，杼中情而属诗。\par
夜炯炯而不寐兮，怀隐忧而历兹。\par
心郁郁而无告兮，众孰可与深谋！\par
欿愁悴而委惰兮，老冉冉而逮之。\par
居处愁以隐约兮，志沉抑而不扬。\par
道壅塞而不通兮，江河广而无梁。\par
愿至昆仑之悬圃兮，采锺山之玉英。\par
揽瑶木之橝枝兮，望阆风之板桐。\par
弱水汩其为难兮，路中断而不通。\par
势不能凌波以径度兮，又无羽翼而高翔。\par
然隐悯而不达兮，独徙倚而彷徉。\par
怅惝罔以永思兮，心纡轸而增伤。\par
倚踌躇以淹留兮，日饥馑而绝粮。\par
廓抱景而独倚兮，超永思乎故乡。\par
廓落寂而无友兮，谁可与玩此遗芳？\par
白日晼晼其將入兮，哀余寿之弗将。\par
车既弊而马罢兮，蹇邅徊而不能行。\par
身既不容于浊世兮，不知进退之宜当。\par
冠崔嵬而切云兮，剑淋离而从横。\par
衣摄叶以储与兮，左袪挂于榑桑\par
右衽拂于不周兮，六合不足以肆行。\par
上同凿枘于伏戏兮，下合矩矱于虞唐。\par
原尊节而式高兮，志犹卑夫禹汤。\par
虽知困其不改操兮，终不以邪枉害方。\par
世并举而好朋兮，壹斗斛而相量。\par
众比周以肩迫兮，贤者远而隐藏。\par
为凤皇作鹑笼兮，虽翕翅其不容。\par
灵皇其不寤知兮，焉陈词而效忠。\par
俗嫉妒而蔽贤兮，孰知余之从容？\par
愿舒志而抽冯兮，庸讵知其吉凶？\par
璋珪杂于甑窐兮，陇廉与孟娵同宫。\par
举世以为恆俗兮，固将愁苦而终穷。\par
幽独转而不寐兮，惟烦懑而盈匈。\par
魂眇眇而驰骋兮，心烦冤之忡忡。\par
志欿憾而不憺兮，路幽昧而甚难。\par
塊独守此曲隅兮，然欿切而永叹。\par
愁修夜而宛转兮，气涫沸其若波。\par
握剞劂而不用兮，操规矩而无所施。\par
骋骐骥于中庭兮，焉能极夫远道？\par
置援狖于棂槛兮，夫何以责其捷巧？\par
驷跛鳖而上山兮，吾固知其不能陞。\par
释管晏而任臧获兮，何权衡之能称？\par
箟簬杂于黀蒸兮，机蓬矢以射革。\par
负檐荷以丈尺兮，欲伸要而不可得。\par
外迫胁于机臂兮，上牵联于矰隿。\par
肩倾侧而不容兮，固陿腹而不得息。\par
务光自投于深渊兮，不获世之尘垢。\par
孰魁摧之可久兮，愿退身而穷处。\par
凿山楹而为室兮，下被衣于水渚。\par
雾露濛濛其晨降兮，云依斐而承宇。\par
虹霓纷其朝霞兮，夕淫淫而淋雨。\par
怊茫茫而无归兮，怅远望此旷野。\par
下垂钓于溪谷兮，上要求于仙者。\par
与赤松而结友兮，比王侨而为耦。\par
使枭杨先导兮，白虎为之前後。\par
浮云雾而入冥兮，骑白鹿而容与。\par
魂眐眐以寄独兮，汨徂往而不归。\par
处卓卓而日远兮，志浩荡而伤怀。\par
鸾凤翔于苍云兮，故矰缴而不能加。\par
蛟龙潜于旋渊兮，身不挂于罔罗。\par
知贪饵而近死兮，不如下游乎清波。\par
宁幽隐以远祸兮，孰侵辱之可为。\par
子胥死而成义兮，屈原沉于汨罗。\par
虽体解其不变兮，岂忠信之可化。\par
志怦怦而内直兮，履绳墨而不颇。\par
执权衡而无私兮，称轻重而不差。\par
摡尘垢之枉攘兮，除秽累而反真。\par
形体白而质素兮，中皎洁而淑清。\par
时猒饫而不用兮，且隐伏而远身。\par
聊窜端而匿迹兮，嗼寂默而无声。\par
独便悁而烦毒兮，焉发愤而筊抒。\par
时暧暧其将罢兮，遂闷叹而无名。\par
伯夷死于首阳兮，卒夭隐而不荣。\par
太公不遇文王兮，身至死而不得逞。\par
怀瑶象而佩琼兮，愿陈列而无正。\par
生天坠之若过兮，忽烂漫而无成。\par
邪气袭余之形体兮，疾憯怛而萌生。\par
原壹见阳春之白日兮，恐不终乎永年。\par

\section{匡机}
\textbf{九怀　王褒}\par\vspace{0.5em}
极运兮不中，来将屈兮困穷。\par
余深愍兮惨怛，愿一列兮无从。\par
乘日月兮上征，顾游心兮鄗酆。\par
弥览兮九隅，彷徨兮兰宫。\par
芷闾兮药房，奋摇兮众芳。\par
菌阁兮蕙楼，观道兮从横。\par
宝金兮委积，美玉兮盈堂。\par
桂水兮潺湲，扬流兮洋洋。\par
蓍蔡兮踊跃，孔鹤兮回翔。\par
抚槛兮远望，念君兮不忘。\par
怫郁兮莫陈，永怀兮内伤。\par

\section{通路}
\textbf{九怀　王褒}\par\vspace{0.5em}
天门兮墬户，孰由兮贤者？\par
无正兮溷厕，怀德兮何睹？\par
假寐兮愍斯，谁可与兮寤语？\par
痛凤兮远逝，畜鴳兮近处。\par
鲸鱏兮幽潜，从虾兮游陼。\par
乘虬兮登阳，载象兮上行。\par
朝发兮葱岭，夕至兮明光。\par
北饮兮飞泉，南采兮芝英。\par
宣游兮列宿，顺极兮彷徉。\par
红采兮骍衣，翠缥兮为裳。\par
舒佩兮綝纚，竦余剑兮干将。\par
腾蛇兮后从，飞駏兮步旁。\par
微观兮玄圃，览察兮瑶光。\par
启匮兮探筴，悲命兮相当。\par
纫蕙兮永辞，将离兮所思。\par
浮云兮容与，道余兮何之？\par
远望兮仟眠，闻雷兮阗阗。\par
阴忧兮感余，惆怅兮自怜。\par

\section{危俊}
\textbf{九怀　王褒}\par\vspace{0.5em}
林不容兮鸣蜩，余何留兮中州？\par
陶嘉月兮总驾，搴玉英兮自修。\par
结荣茝兮逶逝，将去烝兮远游。\par
径岱土兮魏阙，历九曲兮牵牛。\par
聊假日兮相佯，遗光燿兮周流。\par
望太一兮淹息，纡余辔兮自休。\par
晞白日兮皎皎，弥远路兮悠悠。\par
顾列孛兮缥缥，观幽云兮陈浮。\par
钜宝迁兮砏磤，雉咸雊兮相求。\par
泱莽莽兮究志，惧吾心兮懤懤。\par
步余马兮飞柱，览可与兮匹俦。\par
卒莫有兮纤介，永余思兮怞怞。\par

\section{昭世}
\textbf{九怀　王褒}\par\vspace{0.5em}
世溷兮冥昏，违君兮归真。\par
乘龙兮偃蹇，高回翔兮上臻。\par
袭英衣兮缇[糹習]，披华裳兮芳芬。\par
登羊角兮扶舆，浮云漠兮自娱。\par
握神精兮雍容，与神人兮相胥。\par
流星坠兮成雨，进瞵盼兮上丘墟。\par
览旧邦兮滃郁，余安能兮久居。\par
志怀逝兮心懰栗，纡余辔兮踌躇。\par
闻素女兮微歌，听王后兮吹竽。\par
魂悽怆兮感哀，肠回回兮盘纡。\par
抚余佩兮缤纷，高太息兮自怜。\par
使祝融兮先行，令昭明兮开门。\par
驰六蛟兮上征，竦余驾兮入冥。\par
历九州兮索合，谁可与兮终生。\par
忽反顾兮西囿，睹轸丘兮崎倾。\par
横垂涕兮泫流，悲余后兮失灵。\par

\section{尊嘉}
\textbf{九怀　王褒}\par\vspace{0.5em}
季春兮阳阳，列草兮成行。\par
余悲兮兰生，委积兮从横。\par
江离兮遗捐，辛夷兮挤臧。\par
伊思兮往古，亦多兮遭殃。\par
伍胥兮浮江，屈子兮沉湘。\par
运余兮念兹，心内兮怀伤。\par
望淮兮沛沛，滨流兮则逝。\par
榜舫兮下流，东注兮磕磕。\par
蛟龙兮导引，文鱼兮上濑。\par
抽蒲兮陈坐，援芙蕖兮为盖。\par
水跃兮余旌，继以兮微蔡。\par
云旗兮电骛，倏忽兮容裔。\par
河伯兮开门，迎余兮欢欣。\par
顾念兮旧都，怀恨兮艰难。\par
窃哀兮浮萍，汎淫兮无根。\par

\section{蓄英}
\textbf{九怀　王褒}\par\vspace{0.5em}
秋风兮萧萧，舒芳兮振条。\par
微霜兮眇眇，病殀兮鸣蜩。\par
玄鸟兮辞归，飞翔兮灵丘。\par
望谿谷兮滃郁，熊罴兮呴嗥。\par
唐虞兮不存，何故兮久留？\par
临渊兮汪洋，顾林兮忽荒。\par
修余兮袿衣，骑霓兮南上。\par
乘云兮回回，亹亹兮自强。\par
将息兮兰皋，失志兮悠悠。\par
蒶蕴兮霉黧，思君兮无聊。\par
身去兮意存，怆恨兮怀愁。\par

\section{思忠}
\textbf{九怀　王褒}\par\vspace{0.5em}
登九灵兮游神，静女歌兮微晨。\par
悲皇丘兮积葛，众体错兮交纷。\par
贞枝抑兮枯槁，枉车登兮庆云。\par
感余志兮惨栗，心怆怆兮自怜。\par
驾玄螭兮北征，曏吾路兮葱岭。\par
连五宿兮建旄，扬氛气兮为旌。\par
历广漠兮驰骛，览中国兮冥冥。\par
玄武步兮水母，与吾期兮南荣。\par
登华盖兮乘阳，聊逍遥兮播光。\par
抽库娄兮酌醴，援瓟瓜兮接粮。\par
毕休息兮远逝，发玉軔兮西行。\par
惟时俗兮疾正，弗可久兮此方。\par
寤辟摽兮永思，心怫郁兮内伤。\par

\section{陶壅}
\textbf{九怀　王褒}\par\vspace{0.5em}
览杳杳兮世惟，余惆怅兮何归。\par
伤时俗兮溷乱，将奋翼兮高飞。\par
驾八龙兮连蜷，建虹旌兮威夷。\par
观中宇兮浩浩，纷翼翼兮上跻。\par
浮溺水兮舒光，淹低佪兮京沶。\par
屯余车兮索友，睹皇公兮问师。\par
道莫贵兮归真，羡余术兮可夷。\par
吾乃逝兮南娭，道幽路兮九疑。\par
越炎火兮万里，过万首兮嶷嶷。\par
济江海兮蝉蜕，绝北梁兮永辞。\par
浮云郁兮昼昏，霾土忽兮塺々。\par
息阳城兮广夏，衰色罔兮中怠。\par
意晓阳兮燎寤，乃自诊兮在兹。\par
思尧舜兮袭兴，幸咎繇兮获谋。\par
悲九州兮靡君，抚轼叹兮作诗。\par

\section{株昭}
\textbf{九怀　王褒}\par\vspace{0.5em}
悲哉于嗟兮，心内切磋。\par
款冬而生兮，凋彼叶柯。\par
瓦砾进宝兮，捐弃随和。\par
铅刀厉御兮，顿弃太阿。\par
骥垂两耳兮，中坂蹉跎。\par
蹇驴服驾兮，无用日多。\par
修洁处幽兮，贵宠沙劘。\par
凤皇不翔兮，鹑鴳飞扬。\par
乘虹骖蜺兮，载云变化。\par
焦明开路兮，后属青蛇。\par
步骤桂林兮，超骧卷阿。\par
丘陵翔儛兮，谿谷悲歌。\par
神章灵篇兮，赴曲相和。\par
余私娱兹兮，孰哉复加。\par
还顾世俗兮，坏败罔罗。\par
卷佩将逝兮，涕流滂沲。\par
乱曰：皇门开兮照下土，株秽除兮兰芷睹。\par
四佞放兮後得禹，圣舜摄兮昭尧绪，孰能若兮原为辅。\par

\section{逢纷}
\textbf{九叹　刘向}\par\vspace{0.5em}
伊伯庸之末胄兮，谅皇直之屈原。\par
云余肇祖于高阳兮，惟楚怀之婵连。\par
原生受命于贞节兮，鸿永路有嘉名。\par
齐名字于天地兮，并光明于列星。\par
吸精粹而吐氛浊兮，横邪世而不取容。\par
行叩诚而不阿兮，遂见排而逢谗。\par
后听虚而黜实兮，不吾理而顺情。\par
肠愤悁而含怒兮，志迁蹇而左倾。\par
心戃慌其不我与兮，躬速速其不吾亲。\par
辞灵修而陨志兮，吟泽畔之江滨。\par
椒桂罗以颠覆兮，有竭信而归诚。\par
谗夫蔼蔼而漫著兮，曷其不舒予情？\par
始结言于庙堂兮，信中涂而叛之。\par
怀兰蕙与衡芷兮，行中野而散之。\par
声哀哀而怀高丘兮，心愁愁而思旧邦。\par
愿承闲而自恃兮，径淫曀而道壅。\par
颜霉黧以沮败兮，精越裂而衰耄。\par
裳襜襜而含风兮，衣纳纳而掩露。\par
赴江湘之湍流兮，顺波凑而下降。\par
徐徘徊于山阿兮，飘风来之洶洶。\par
驰余车兮玄石，步余马兮洞庭。\par
平明发兮苍梧，夕投宿兮石城。\par
芙蓉盖而菱华车兮，紫贝阙而玉堂。\par
薜荔饰而陆离荐兮，鱼鳞衣而白蜺裳。\par
登逢龙而下陨兮，违故都之漫漫。\par
思南郢之旧俗兮，肠一夕而九运。\par
扬流波之潢潢兮，体溶溶而东回。\par
心怊怅以永思兮，意晻晻而日颓。\par
白露纷以涂涂兮，秋风浏以萧萧。\par
身永流而不还兮，魂长逝而常愁。\par
叹曰：\par
譬彼流水纷扬磕兮，波逢汹涌濆壅滂兮。\par
揄扬涤荡飘流陨往触崟石兮，龙卬脟圈缭戾宛转阻相薄兮，\par
遭纷逢凶蹇离尤兮，垂文扬采遗将来兮。\par

\section{离世}
\textbf{九叹　刘向}\par\vspace{0.5em}
灵怀其不吾知兮，灵怀其不吾闻。\par
就灵怀之皇祖兮，愬灵怀之鬼神。\par
灵怀曾不吾与兮，即听夫人之谀辞。\par
余辞上参于天坠兮，旁引之于四时。\par
指日月使延照兮，抚招摇以质正。\par
立师旷俾端辞兮，命咎繇使并听。\par
兆出名曰正则兮，卦发字曰灵均。\par
余幼既有此鸿节兮，长愈固而弥纯。\par
不从俗而诐行兮，直躬指而信志。\par
不枉绳以追曲兮，屈情素以从事。\par
端余行其如玉兮，述皇舆之踵迹。\par
群阿容以晦光兮，皇舆覆以幽辟。\par
舆中涂以回畔兮，驷马惊而横奔。\par
执组者不能制兮，必折轭而摧辕。\par
断镳衔以驰骛兮，暮去次而敢止。\par
路荡荡其无人兮，遂不禦乎千里。\par
身衡陷而下沉兮，不可获而复登。\par
不顾身之卑贱兮，惜皇舆之不兴。\par
出国门而端指兮，冀壹寤而锡还。\par
哀仆夫之坎毒兮，屡离忧而逢患。\par
九年之中不吾反兮，思彭咸之水游。\par
惜师延之浮渚兮，赴汨罗之长流。\par
遵江曲之逶移兮，触石碕而衡游。\par
波澧澧而扬浇兮，顺长濑之浊流。\par
凌黄沱而下低兮，思还流而复反。\par
玄舆驰而并集兮，身容与而日远。\par
棹舟杭以横濿兮，济湘流而南极。\par
立江界而长吟兮，愁哀哀而累息。\par
情慌忽以忘归兮，神浮游以高历。\par
心蛩蛩而怀顾兮，魂眷眷而独逝。\par
叹曰：\par
余思旧邦心依违兮，日暮黄昏羌幽悲兮，\par
去郢东迁余谁慕兮，谗夫党旅其以兹故兮，\par
河水淫淫情所愿兮，顾瞻郢路终不返兮。\par

\section{怨思}
\textbf{九叹　刘向}\par\vspace{0.5em}
惟郁郁之忧毒兮，志坎壈而不违。\par
身憔悴而考旦兮，日黄昏而长悲。\par
闵空宇之孤子兮，哀枯杨之冤雏。\par
孤雌吟于高墉兮，鸣鸠栖于桑榆。\par
玄蝯失于潜林兮，独偏弃而远放。\par
征夫劳于周行兮，处妇愤而长望。\par
申诚信而罔违兮，情素洁于纽帛。\par
光明齐于日月兮，文采耀燿于玉石。\par
伤压次而不发兮，思沉抑而不扬。\par
芳懿懿而终败兮，名靡散而不彰。\par
背玉门以奔骛兮，蹇离尤而干诟。\par
若龙逢之沉首兮，王子比干之逢醢。\par
念社稷之几危兮，反为雠而见怨。\par
思国家之离沮兮，躬获愆而结难。\par
若青蝇之伪质兮，晋骊姬之反情。\par
恐登阶之逢殆兮，故退伏于末庭。\par
孽臣之号咷兮，本朝芜而不治。\par
犯颜色而触谏兮，反蒙辜而被疑。\par
菀蘼芜与菌若兮，渐藁本于洿渎。\par
淹芳芷于腐井兮，弃鸡骇于筐簏。\par
执棠谿以刜蓬兮，秉干将以割肉。\par
筐泽泻以豹鞟兮，破荆和以继筑。\par
时溷浊犹未清兮，世殽乱犹未察。\par
欲容与以俟时兮，惧年岁之既晏。\par
顾屈节以从流兮，心巩巩而不夷。\par
宁浮沅而驰骋兮，下江湘以邅回。\par
叹曰：\par
山中槛槛余伤怀兮，征夫皇皇其孰依兮，\par
经营原野杳冥冥兮，乘骐骋骥舒吾情兮，\par
归骸旧邦莫谁语兮，长辞远逝乘湘去兮。\par

\section{远逝}
\textbf{九叹　刘向}\par\vspace{0.5em}
志隐隐而郁怫兮，愁独哀而冤结。\par
肠纷纭以缭转兮，涕渐渐其若屑。\par
情慨慨而长怀兮，信上皇而质正。\par
合五岳与八灵兮，讯九鬿与六神。\par
指列宿以白情兮，诉五帝以置辞。\par
北斗为我折中兮，太一为余听之。\par
云服阴阳之正道兮，御后土之中和。\par
佩苍龙之蚴虬兮，带隐虹之逶蛇。\par
曳彗星之皓旰兮，抚朱爵与鵔鸃。\par
游清灵之飒戾兮，服云衣之披披。\par
杖玉策与朱旗兮，垂明月之玄珠。\par
举霓旌之墆翳兮，建黄纁之总旄。\par
躬纯粹而罔愆兮，承皇考之妙仪。\par
惜往事之不合兮，横汨罗而下沥。\par
乘隆波而南渡兮，逐江湘之顺流。\par
赴阳侯之潢洋兮，下石濑而登洲。\par
陆魁堆以蔽视兮，云冥冥而闇前。\par
山峻高以无垠兮，遂曾闳而迫身。\par
雪雰雰而薄木兮，云霏霏而陨集。\par
阜隘狭而幽险兮，石嵾嵯以翳日。\par
悲故乡而发忿兮，去余邦之弥久。\par
背龙门而入河兮，登大坟而望夏首。\par
横舟航而济湘兮，耳聊啾而戃慌。\par
波淫淫而周流兮，鸿溶溢而滔荡。\par
路曼曼其无端兮，周容容而无识。\par
引日月以指极兮，少须臾而释思。\par
水波远以冥冥兮，眇不睹其东西。\par
顺风波以南北兮，雾宵晦以纷纷。\par
日杳杳以西颓兮，路长远而窘迫。\par
欲酌醴以娱忧兮，蹇骚骚而不释。\par
叹曰：\par
飘风蓬龙埃坲々兮，草木摇落时槁悴兮，\par
遭倾遇祸不可救兮，长吟永欷涕究究兮，\par
舒情陈诗冀以自免兮，颓流下陨身日远兮。\par

\section{惜贤}
\textbf{九叹　刘向}\par\vspace{0.5em}
览屈氏之离骚兮，心哀哀而怫郁。\par
声嗷嗷以寂寥兮，顾仆夫之憔悴。\par
拨谄谀而匡邪兮，切淟涊之流俗。\par
荡渨涹之奸咎兮，夷蠢蠢之溷浊。\par
怀芬香而挟蕙兮，佩江蓠之婓婓。\par
握申椒与杜若兮，冠浮云之峨峨。\par
登长陵而四望兮，览芷圃之蠡蠡。\par
游兰皋与蕙林兮，睨玉石之嵾嵯。\par
扬精华以炫燿兮，芳郁渥而纯美。\par
结桂树之旖旎兮，纫荃蕙与辛夷。\par
芳若兹而不御兮，捐林薄而菀死。\par
驱子侨之奔走兮，申徒狄之赴渊。\par
若由夷之纯美兮，介子推之隐山。\par
晋申生之离殃兮，荆和氏之泣血。\par
吴申胥之抉眼兮，王子比干之横废。\par
欲卑身而下体兮，心隐恻而不置。\par
方圜殊而不合兮，钩绳用而异态。\par
欲俟时于须臾兮，日阴曀其将暮。\par
时迟迟其日进兮，年忽忽而日度。\par
妄周容而入世兮，内距闭而不开。\par
俟时风之清激兮，愈氛雾其如塺。\par
进雄鸠之耿耿兮，谗介介而蔽之。\par
默顺风以偃仰兮，尚由由而进之。\par
心懭悢以冤结兮，情舛错以曼忧。\par
搴薜荔于山野兮，采撚支于中洲。\par
望高丘而叹涕兮，悲吸吸而长怀。\par
孰契契而委栋兮，日晻晻而下颓。\par
叹曰：\par
江湘油油长流汩兮，挑揄扬汰荡迅疾兮。\par
忧心展转愁怫郁兮，冤结未舒长隐忿兮，\par
丁时逢殃可奈何兮，劳心悁悁涕滂沱兮。\par

\section{忧苦}
\textbf{九叹　刘向}\par\vspace{0.5em}
悲余心之悁悁兮，哀故邦之逢殃。\par
辞九年而不复兮，独茕茕而南行。\par
思余俗之流风兮，心纷错而不受。\par
遵野莽以呼风兮，步从容于山廋。\par
巡陆夷之曲衍兮，幽空虚以寂寞。\par
倚石岩以流涕兮，忧憔悴而无乐。\par
登巑岏以长企兮，望南郢而闚之。\par
山修远其辽辽兮，涂漫漫其无时。\par
听玄鹤之晨鸣兮，于高冈之峨峨。\par
独愤积而哀娱兮，翔江洲而安歌。\par
三鸟飞以自南兮，览其志而欲北。\par
原寄言于三鸟兮，去飘疾而不可得。\par
欲迁志而改操兮，心纷结其未离。\par
外彷徨而游览兮，内恻隐而含哀。\par
聊须臾以时忘兮，心渐渐其烦错。\par
原假簧以舒忧兮，志纡郁其难释。\par
叹《离骚》以扬意兮，犹未殫于《九章》。\par
长嘘吸以于悒兮，涕横集而成行。\par
伤明珠之赴泥兮，鱼眼玑之坚藏。\par
同驽骡与乘駔兮，杂斑駮与阘茸。\par
葛藟虆于桂树兮，鸱鸮集于木兰。\par
偓促谈于廊庙兮，律魁放乎山间。\par
恶虞氏之箫《韶》兮，好遗风之《激楚》。\par
潜周鼎于江淮兮，爨土鬵于中宇。\par
且人心之持旧兮，而不可保长。\par
邅彼南道兮，征夫宵行。\par
思念郢路兮，还顾睠睠。\par
涕流交集兮，泣下涟涟。\par
叹曰：\par
登山长望中心悲兮，菀彼青青泣如颓兮，\par
留思北顾涕渐渐兮，折锐摧矜凝氾滥兮，\par
念我茕茕魂谁求兮，仆夫慌悴散若流兮。\par

\section{愍命}
\textbf{九叹　刘向}\par\vspace{0.5em}
昔皇考之嘉志兮，喜登能而亮贤。\par
情纯洁而罔薉兮，姿盛质而无愆。\par
放佞人与谄谀兮，斥谗夫与便嬖。\par
亲忠正之悃诚兮，招贞良与明智。\par
心溶溶其不可量兮，情澹澹其若渊。\par
回邪辟而不能入兮，诚原藏而不可迁。\par
逐下袟于後堂兮，迎虙妃于伊雒。\par
刜谗贼于中廇兮，选吕管于榛薄。\par
丛林之下无怨士兮，江河之畔无隐夫。\par
三苗之徒以放逐兮，伊皋之伦以充庐。\par
今反表以为里兮，颠裳以为衣。\par
戚宋万于两楹兮，废周邵于遐夷。\par
却骐骥以转运兮，腾驴骡以驰逐。\par
蔡女黜而出帷兮，戎妇入而綵绣服。\par
庆忌囚于阱室兮，陈不占战而赴围。\par
破伯牙之号钟兮，挟人筝而弹纬。\par
藏瑉石于金匮兮，捐赤瑾于中庭。\par
韩信蒙于介胄兮，行夫将而攻城。\par
莞芎弃于泽洲兮，瓟瓥蠹于筐簏。\par
麒麟奔于九皋兮，熊罴群而逸囿。\par
折芳枝与琼华兮，树枳棘与薪柴。\par
掘荃蕙与射干兮，耘藜藿与蘘荷。\par
惜今世其何殊兮，远近思而不同。\par
或沉沦其无所达兮，或清激其无所通。\par
哀余生之不当兮，独蒙毒而逢尤。\par
虽謇謇以申志兮，君乖差而屏之。\par
诚惜芳之菲菲兮，反以兹为腐也。\par
怀椒聊之蔎蔎兮，乃逢纷以罹诟也。\par
叹曰：\par
嘉皇既殁终不返兮，山中幽险郢路远兮。\par
谗人諓諓孰可愬兮，征夫罔极谁可语兮。\par
行吟累欷声喟喟兮，怀忧含戚何侘傺兮。\par

\section{思古}
\textbf{九叹　刘向}\par\vspace{0.5em}
冥冥深林兮，树木郁郁。\par
山参差以崭岩兮，阜杳杳以蔽日。\par
悲余心之悁悁兮，目眇眇而遗泣。\par
风骚屑以摇木兮，云吸吸以湫戾。\par
悲余生之无欢兮，愁倥傯于山陆。\par
旦徘徊于长阪兮，夕彷徨而独宿。\par
发披披以鬤鬤兮，躬劬劳而瘏悴。\par
魂俇俇而南行兮，泣霑襟而濡袂。\par
心婵媛而无告兮，口噤闭而不言。\par
违郢都之旧闾兮，回湘、沅而远迁。\par
念余邦之横陷兮，宗鬼神之无次。\par
闵先嗣之中绝兮，心惶惑而自悲。\par
聊浮游于山陿兮，步周流于江畔。\par
临深水而长啸兮，且倘佯而氾观。\par
兴离骚之微文兮，冀灵修之壹悟。\par
还余车于南郢兮，复往轨于初古。\par
道修远其难迁兮，伤余心之不能已。\par
背三五之典刑兮，绝洪范之辟纪。\par
播规矩以背度兮，错权衡而任意。\par
操绳墨而放弃兮，倾容幸而侍侧。\par
甘棠枯于丰草兮，藜棘树于中庭。\par
西施斥于北宫兮，仳倠倚于弥楹。\par
乌获戚而骖乘兮，燕公操于马圉。\par
蒯聩登于清府兮，咎繇弃而在野。\par
盖见兹以永叹兮，欲登阶而狐疑。\par
乘白水而高骛兮，因徙弛而长词。\par
叹曰：\par
倘佯垆阪沼水深兮，容与汉渚涕淫淫兮，锺牙已死谁为声兮？\par
纤阿不御焉舒情兮，曾哀悽欷心离离兮，还顾高丘泣如洒兮。\par

\section{远游}
\textbf{九叹　刘向}\par\vspace{0.5em}
悲余性之不可改兮，屡惩艾而不迻。\par
服觉皓以殊俗兮，貌揭揭以巍巍。\par
譬若王侨之乘云兮，载赤霄而凌太清。\par
欲与天地参寿兮，与日月而比荣。\par
登昆仑而北首兮，悉灵圉而来谒。\par
选鬼神于太阴兮，登阊阖于玄阙。\par
回朕车俾西引兮，褰虹旗于玉门。\par
驰六龙于三危兮，朝西灵于九滨。\par
结余轸于西山兮，横飞谷以南征。\par
绝都广以直指兮，历祝融于硃冥。\par
枉玉衡于炎火兮，委两馆于咸唐。\par
贯澒濛以东朅兮，维六龙于扶桑。\par
周流览于四海兮，志升降以高驰。\par
徵九神于回极兮，建虹采以招指。\par
驾鸾凤以上游兮，从玄鹤与鹪明。\par
孔鸟飞而送迎兮，腾群鹤于瑶光。\par
排帝宫与罗囿兮，升县圃以眩灭。\par
结琼枝以杂佩兮，立长庚以继日。\par
凌惊雷以轶骇电兮，缀鬼谷于北辰。\par
鞭风伯使先驱兮，囚灵玄于虞渊。\par
遡高风以低佪兮，览周流于朔方。\par
就颛顼而敶辞兮，考玄冥于空桑。\par
旋车逝于崇山兮，奏虞舜于苍梧。\par
济杨舟于会稽兮，就申胥于五湖。\par
见南郢之流风兮，殒余躬于沅湘。\par
望旧邦之黯黮兮，时溷浊其犹未央。\par
怀兰茝之芬芳兮，妒被离而折之。\par
张绛帷以襜襜兮，风邑邑而蔽之。\par
日暾暾其西舍兮，阳焱焱而复顾。\par
聊假日以须臾兮，何骚骚而自故。\par
叹曰：\par
譬彼蛟龙乘云浮兮，汎淫澒溶纷若雾兮。\par
潺湲轇轕雷动电发馺高举兮。\par
升虚凌冥沛浊浮清入帝宫兮，摇翘奋羽驰风骋雨游无穷兮。\par

\section{逢尤}
\textbf{九思　王逸}\par\vspace{0.5em}
悲兮愁，哀兮忧！\par
天生我兮当闇时，被诼谮兮虚获尤。\par
心烦憒兮意无聊，严载驾兮出戏游。\par
周八极兮历九州，求轩辕兮索重华。\par
世既卓兮远眇眇，握佩玖兮中路躇。\par
羡咎繇兮建典谟，懿风后兮受瑞图。\par
愍余命兮遭六极，委玉质兮於泥涂。\par
遽傽遑兮驱林泽，步屏营兮行丘阿。\par
车軏折兮马虺颓，惷怅立兮涕滂沱。\par
思丁文兮圣明哲，哀平差兮迷谬愚。\par
吕傅举兮殷周兴，忌嚭专兮郢吴虚。\par
仰长叹兮气噎结，悒殟绝兮咶复苏。\par
虎兕争兮於廷中，豺狼斗兮我之隅。\par
云雾会兮日冥晦，飘风起兮扬尘埃。\par
走鬯罔兮乍东西，欲窜伏兮其焉如？\par
念灵闺兮隩重深，原竭节兮隔无由。\par
望旧邦兮路逶随，忧心悄兮志勤劬。\par
魂茕茕兮不遑寐，目眽眽兮寤终朝。\par

\section{怨上}
\textbf{九思　王逸}\par\vspace{0.5em}
令尹兮謷謷，群司兮譨譨。\par
哀哉兮淈淈，上下兮同流。\par
菽藟兮蔓衍，芳虈兮挫枯。\par
硃紫兮杂乱，曾莫兮别诸。\par
倚此兮岩穴，永思兮窈悠。\par
嗟怀兮眩惑，用志兮不昭。\par
将丧兮玉斗，遗失兮钮枢。\par
我心兮煎熬，惟是兮用忧。\par
进恶兮九旬，复顾兮彭务。\par
拟斯兮二踪，未知兮所投。\par
谣吟兮中野，上察兮璇玑。\par
大火兮西睨，摄提兮运低。\par
雷霆兮硠磕，雹霰兮霏霏。\par
奔电兮光晃，凉风兮怆悽。\par
鸟兽兮惊骇，相从兮宿栖。\par
鸳鸯兮噰噰，狐狸兮徾徾。\par
哀吾兮介特，独处兮罔依。\par
蝼蛄兮鸣东，蟊蠽兮号西。\par
蛓缘兮我裳，蠋入兮我怀。\par
虫豸兮夹余，惆怅兮自悲。\par
伫立兮忉怛，心结縎兮折摧。\par

\section{疾世}
\textbf{九思　王逸}\par\vspace{0.5em}
周徘徊兮汉渚，求水神兮灵女。\par
嗟此国兮无良，媒女诎兮謰謱。\par
鴳雀列兮譁讙，鸲鹆鸣兮聒余。\par
抱昭华兮宝璋，欲衒鬻兮莫取。\par
言旋迈兮北徂，叫我友兮配耦。\par
日阴曀兮未光，阒睄窕兮靡睹。\par
纷载驱兮高驰，将谘询兮皇羲。\par
遵河皋兮周流，路变易兮时乖。\par
濿沧海兮东游，沐盥浴兮天池。\par
访太昊兮道要，云靡贵兮仁义。\par
志欣乐兮反征，就周文兮邠歧。\par
秉玉英兮结誓，日欲暮兮心悲。\par
惟天禄兮不再，背我信兮自违。\par
逾陇堆兮渡漠，过桂车兮合黎。\par
赴昆山兮馽騄，从邛遨兮栖迟。\par
吮玉液兮止渴，齧芝华兮疗饥。\par
居嵺廓兮尠畴，远梁昌兮几迷。\par
望江汉兮濩渃，心紧絭兮伤怀。\par
时昢昢兮且旦，尘莫莫兮未晞。\par
忧不暇兮寝食，吒增叹兮如雷。\par

\section{悯上}
\textbf{九思　王逸}\par\vspace{0.5em}
哀世兮睩睩，諓諓兮嗌喔。\par
众多兮阿媚，骪靡兮成俗。\par
贪枉兮党比，贞良兮茕独。\par
鹄窜兮枳棘，鹈集兮帷幄。\par
罽蕠兮青葱，槁本兮萎落。\par
睹斯兮伪惑，心为兮隔错。\par
逡巡兮圃薮，率彼兮畛陌。\par
川谷兮渊渊，山阜兮峉峉。\par
丛林兮崟崟，株榛兮岳岳。\par
霜雪兮漼溰，冰冻兮洛泽。\par
东西兮南北，罔所兮归薄。\par
庇廕兮枯树，匍匐兮岩石。\par
蜷跼兮寒局数，独处兮志不申。\par
年齿尽兮命迫促，魁垒挤摧兮常困辱。\par
含忧强老兮愁无乐，须发苎悴兮顠鬓白。\par
思灵泽兮一膏沐，怀兰英兮把琼若，待天明兮立踯躅。\par
云蒙蒙兮电鯈烁，孤雌惊兮鸣呴呴。\par
思怫郁兮肝切剥，忿悁悒兮孰诉告。\par

\section{遭厄}
\textbf{九思　王逸}\par\vspace{0.5em}
悼屈子兮遭厄，沉王躬兮湘汨。\par
何楚国兮难化，迄于今兮不易。\par
士莫志兮羔裘，竞佞谀兮谗阋。\par
指正义兮为曲，訿玉璧兮为石。\par
殦雕游兮华屋，鵔鸃栖兮柴蔟。\par
起奋迅兮奔走，违群小兮謑訽。\par
载青云兮上昇，適昭明兮所处。\par
蹑天衢兮长驱，踵九阳兮戏荡。\par
越云汉兮南济，秣余马兮河鼓。\par
云霓纷兮晻翳，参辰回兮颠倒。\par
逢流星兮问路，顾我指兮从左。\par
俓娵觜兮直驰，御者迷兮失轨。\par
遂踢达兮邪造，与日月兮殊道。\par
志阏绝兮安如，哀所求兮不耦。\par
攀天阶兮下视，见鄢郢兮旧宇。\par
意逍遥兮欲归，众秽盛兮沓沓。\par
思哽饐兮诘诎，涕流澜兮如雨。\par

\section{悼乱}
\textbf{九思　王逸}\par\vspace{0.5em}
嗟嗟兮悲夫，殽乱兮纷挐。\par
茅丝兮同综，冠屦兮共絇。\par
督万兮侍宴，周邵兮负刍。\par
白龙兮见射，灵龟兮执拘。\par
仲尼兮困厄，邹衍兮幽囚。\par
伊余兮念兹，奔遁兮隐居。\par
将升兮高山，上有兮猴猿。\par
欲入兮深谷，下有兮虺蛇。\par
左见兮鸣鵙，右睹兮呼枭。\par
惶悸兮失气，踊跃兮距跳。\par
便旋兮中原，仰天兮增叹。\par
菅蒯兮楙莽，雚苇兮仟眠。\par
鹿蹊兮躖躖，貒貉兮蟫蟫。\par
鹯鹞兮轩轩，鹑鹌兮甄甄。\par
哀我兮寡独，靡有兮齐伦。\par
意欲兮沉吟，迫日兮黄昏。\par
玄鹤兮高飞，曾逝兮青冥。\par
鶬鶊兮喈喈，山鹊兮嚶嚶。\par
鸿鸬兮振翅，归雁兮于征。\par
吾志兮觉悟，怀我兮圣京。\par
垂屣兮将起，跓俟兮硕明。\par

\section{伤时}
\textbf{九思　王逸}\par\vspace{0.5em}
惟昊天兮昭灵，阳气发兮清明。\par
风習習兮和暖，百草萌兮华荣。\par
堇荼茂兮扶疏，蘅芷彫兮莹嫇。\par
愍贞良兮遇害，将夭折兮碎糜。\par
时混混兮浇饡，哀当世兮莫知。\par
览往昔兮俊彦，亦诎辱兮系纍。\par
管束缚兮桎梏，百贸易兮传卖。\par
遭桓缪兮识举，才德用兮列施。\par
且从容兮自慰，玩琴书兮游戏。\par
迫中国兮迮陿，吾欲之兮九夷。\par
超五岭兮嵯峨，观浮石兮崔嵬。\par
陟丹山兮炎野，屯余车兮黄支。\par
就祝融兮稽疑，嘉己行兮无为。\par
乃回朅兮北逝，遇神孈兮宴娭。\par
欲静居兮自娱，心愁慼兮不能。\par
放余辔兮策驷，忽飙腾兮浮云。\par
蹠飞杭兮越海，从安期兮蓬莱。\par
缘天梯兮北上，登太一兮玉台。\par
使素女兮鼓簧，乘戈和兮讴谣。\par
声噭誂兮清和，音晏衍兮要婬。\par
咸欣欣兮酣乐，余眷眷兮独悲。\par
顾章华兮太息，志恋恋兮依依。\par

\section{哀岁}
\textbf{九思　王逸}\par\vspace{0.5em}
旻天兮清凉，玄气兮高朗。\par
北风兮潦洌，草木兮苍唐。\par
蛜蚗兮噍噍，蝍蛆兮穰穰。\par
岁忽忽兮惟暮，余感时兮悽怆。\par
伤俗兮泥浊，矇蔽兮不章。\par
宝彼兮沙砾，捐此兮夜光。\par
椒瑛兮湟汙，葈耳兮充房。\par
摄衣兮缓带，操我兮墨阳。\par
昇车兮命仆，将驰兮四荒。\par
下堂兮见虿，出门兮触螽。\par
巷有兮蚰蜓，邑多兮螳螂。\par
睹斯兮嫉贼，心为兮切伤。\par
俛念兮子胥，仰怜兮比干。\par
投剑兮脱冕，龙屈兮蜿蟤。\par
潜藏兮山泽，匍匐兮丛攒。\par
窥见兮溪涧，流水兮沄沄。\par
鼋鼍兮欣欣，鱣鲇兮延延。\par
群行兮上下，骈罗兮列陈。\par
自恨兮无友，特处兮茕茕。\par
冬夜兮陶陶，雨雪兮冥冥。\par
神光兮颎颎，鬼火兮荧荧。\par
修德兮困控，愁不聊兮遑生。\par
忧纡兮郁郁，恶所兮写情。\par

\section{守志}
\textbf{九思　王逸}\par\vspace{0.5em}
陟玉峦兮逍遥，览高冈兮峣峣。\par
桂树列兮纷敷，吐紫华兮布条。\par
实孔鸾兮所居，今其集兮惟鸮。\par
乌鹊惊兮哑哑，余顾盼兮怊怊。\par
彼日月兮闇昧，障覆天兮祲氛。\par
伊我后兮不聪，焉陈诚兮效忠。\par
摅羽翮兮超俗，游陶遨兮养神。\par
乘六蛟兮蜿蝉，遂驰骋兮升云。\par
扬彗光兮为旗，秉电策兮为鞭。\par
朝晨发兮鄢郢，食时至兮增泉。\par
绕曲阿兮北次，造我车兮南端。\par
谒玄黄兮纳贽，崇忠贞兮弥坚。\par
历九宫兮遍观，睹秘藏兮宝珍。\par
就传说兮骑龙，与织女兮合婚。\par
举天罼兮掩邪，彀天弧兮射奸。\par
随真人兮翱翔，食元气兮长存。\par
望太微兮穆穆，睨三阶兮炳分。\par
相辅政兮成化，建烈业兮垂勋。\par
目瞥瞥兮西没，道遐回兮阻叹。\par
志稸积兮未通，怅敞罔兮自怜。\par
乱曰：\par
天庭明兮云霓藏，三光朗兮镜万方。\par
斥蜥蜴兮进龟龙，策谋从兮翼机衡。\par
配稷契兮恢唐功，嗟英俊兮未为双。\par

\end{document}